\chapter{能控性与能观性：结构、对偶与最小实现}
\label{ch:controllability-obsv}

% ════════════════════════════════════════════════════════════════
%  C0 现代控制理论基础子部 · 第 C0-2 章
%  主源：刘豹《现代控制理论》ch3（孤立方块直觉·卡尔曼分解·标准型·最小实现·§3.10 零极点对消）
%  辅源/互校：张嗣瀛-高立群《现代控制理论》ch4（标准形最干净·时变 Gramian·输出能控·PBH·最小实现）
%  密度梯度：导论 ch:control 已"给结论会用"，本章"讲透+推一遍+建直觉"，严格证明/前沿交 A–Q。
%  教学层遵循 v5.0（动机→反面→历史→理论→陷阱→练习 + 符号表/速查/故障排查）；图全部 TikZ 内联重画。
% ════════════════════════════════════════════════════════════════

\begin{practice}[前置自测]
开始之前，先试答下列问题。若已能流畅作答，可只速读本章表格与速查卡；若卡住，正是本章要为你解决的。
\begin{enumerate}
  \item 系统 $\dot{\mathbf{x}}=\mathbf{A}\mathbf{x}+\mathbf{B}\mathbf{u}$ 的“可控”到底是说什么能做到、什么做不到？它和“稳定”是一回事吗？
  \item 导论 \cref{eq:ctrl-cmat} 给出可控性矩阵 $\mathcal{C}=[\mathbf{B}\ \mathbf{A}\mathbf{B}\ \cdots\ \mathbf{A}^{n-1}\mathbf{B}]$ 满秩即可控。\emph{为什么}只算到 $\mathbf{A}^{n-1}$ 就够、不必算 $\mathbf{A}^n,\mathbf{A}^{n+1},\dots$？
  \item 把系统对角化后，$\dot{x}_i=\lambda_i x_i+b_i u$。若某个 $b_i=0$，这一维状态会怎样？它“跑去哪了”、为什么控制不到？
  \item 可控性与可观性，一个讲“输入能否驱动状态”、一个讲“输出能否反推状态”——它们之间有没有一条精确的数学桥，让你证一个就白得另一个？
  \item 一个传递函数 $W(s)=\dfrac{s+2.5}{(s+2.5)(s-1)}$ 分子分母约去了 $(s+2.5)$。被约掉的那个极点“消失”了吗？它在状态空间里以什么形式存在、为什么危险？
  \item 同一个传递函数可以对应无穷多个状态空间实现，维数各不相同。其中“最小”的那个有什么特殊？它和可控、可观是什么关系？
\end{enumerate}
\end{practice}

\section*{本章目标}
学完本章，你应当能够：
\begin{enumerate}
  \item \textbf{说清}可控性 / 可达性 / 可观性的定义（连续、离散、时变三种语境），并辨析“某状态可控”与“系统状态完全可控”的区别；
  \item \textbf{建立}可控可观的\emph{结构直觉}——用约旦型“孤立方块”一眼看出哪些模态与输入 / 输出脱钩，理解秩判据背后“有没有与输入无关的孤立部分”这一物理图像；
  \item \textbf{独立推导} Kalman 秩判据（连续与离散），并解释 Cayley--Hamilton 定理为何把“无穷步可达”压缩为“$n$ 步可达”；
  \item \textbf{掌握}时变系统的 Gram 矩阵判据，理解它为何是定常 $(\mathbf{A},\mathbf{B})$ 秩判据的“一般形式”；
  \item \textbf{运用}可控--可观\emph{对偶原理}，把可控性结论一键翻译为可观性结论（含时变情形的转置逆关系）；
  \item \textbf{推导}并使用可控 / 可观\textbf{标准型}（友矩阵形），知道它们为什么是状态反馈 / 观测器设计的天然坐标；
  \item \textbf{完成} Kalman 结构分解——把任意 LTI 系统拆成“可控可观 / 可控不可观 / 不可控可观 / 不可控不可观”四块，并理解传递函数为何只“看得见”可控且可观的那一块；
  \item \textbf{求取}传递函数的\textbf{最小实现}，并掌握单输入单输出情形“最小实现 $\iff$ 可控且可观 $\iff$ 无零极点对消”这条贯通定理；
  \item \textbf{把握}本章在全 P4 中的地位：可控 / 可观是 LQR、观测器、$H_\infty$、CBF 设计的\emph{前提条件}，标准型与最小实现是它们的 ground-truth。
\end{enumerate}

\subsection*{与导论的分工}
本章是 P4 控制部「现代控制理论基础」子部（C0）的第二章，紧接 \cref{ch:statespace-basics}（状态空间建模与求解）之后。\textbf{它与导论 \cref{ch:control} 的关系是“同一主题、不同密度”}：导论 \cref{sec:ctrl-statespace} 已把可控性矩阵（\cref{eq:ctrl-cmat}）、Gramian（\cref{eq:ctrl-gramian}）、PBH 判据（\cref{eq:ctrl-pbh}）、可观性矩阵（\cref{eq:ctrl-omat}）、对偶定理（\cref{thm:ctrl-duality}）以\emph{导论密度}陈述清楚——“给结论、会用”。本章不重复这些编号，而是\textbf{把它们讲透：先建物理直觉与几何图像，再把每条判据推一遍，最后铺到标准型、结构分解、最小实现这一整套“结构理论”}。一句话：导论让你\emph{会用}可控可观判据，本章让你\emph{会推、且理解为什么}。凡导论已立的 def / thm，本章开篇半句交叉引用接住即下沉展开（如“展开 \cref{thm:ctrl-gramian-rank}”），决不另起炉灶。

\subsection*{本章知识导航}
本章是一条“\emph{从一个状态可控 / 可观的判别，走到整个系统的结构分解与最小实现}”的主线。结构如下：

\begin{center}
\small
\begin{tikzpicture}[
  node distance=4mm and 7mm, >=Stealth,
  blk/.style={draw, rounded corners, align=center, font=\small, inner sep=3pt, fill=main!6, draw=main!60!black},
  grp/.style={draw, rounded corners, align=center, font=\small, inner sep=3pt, fill=second!6, draw=second!60!black},
  red/.style={draw, rounded corners, align=center, font=\small, inner sep=3pt, fill=third!8, draw=third!60!black},
  ln/.style={->, thick, gray!60!black}]
\node[blk] (def) {定义\\可控/可达/可观};
\node[blk, right=of def] (jordan) {约旦型直觉\\“孤立方块”};
\node[blk, right=of jordan] (rank) {秩判据\\Kalman / PBH};
\node[grp, below=8mm of def] (tv) {时变\\Gram 判据};
\node[grp, right=of tv] (dual) {对偶原理\\$\mathbf{A}\!\to\!\mathbf{A}^\top$};
\node[grp, right=of dual] (canon) {可控/可观\\标准型};
\node[red, below=8mm of tv] (decomp) {Kalman\\结构分解};
\node[red, right=of decomp] (minreal) {最小实现};
\node[red, right=of minreal] (cancel) {零极点对消\\（SISO 贯通）};
\draw[ln] (def)--(jordan); \draw[ln] (jordan)--(rank);
\draw[ln] (rank.south) -- ++(0,-4mm) -| (tv.north);
\draw[ln] (tv)--(dual); \draw[ln] (dual)--(canon);
\draw[ln] (canon.south) -- ++(0,-4mm) -| (decomp.north);
\draw[ln] (decomp)--(minreal); \draw[ln] (minreal)--(cancel);
\end{tikzpicture}
\end{center}

\noindent\textbf{推荐路径}：\cref{sec:co-def,sec:co-jordan} 是任何读者的主干（定义 + 孤立方块直觉，这是全章物理图像的根）；\cref{sec:co-rank,sec:co-observe} 把秩判据推透并接上对偶；\cref{sec:co-timevarying} 的时变 Gram 判据面向\emph{沿轨迹线性化}的机器人系统（LTV），可在需要时回看；\cref{sec:co-canonical}（标准型）与 \cref{sec:co-decomp}（结构分解）是综合 / 最优 / 估计章的 ground-truth；\cref{sec:co-minreal}（最小实现）与 \cref{sec:co-cancel}（零极点对消）把“可控可观”与经典传递函数语言焊接，是理解“现代 vs 经典”的关键一节。

\subsection*{前置知识桥接}
\cref{ch:statespace-basics} 已给出状态空间表达式 $\dot{\mathbf{x}}=\mathbf{A}\mathbf{x}+\mathbf{B}\mathbf{u},\ \mathbf{y}=\mathbf{C}\mathbf{x}+\mathbf{D}\mathbf{u}$、线性变换 $\mathbf{x}=\mathbf{T}\bar{\mathbf{x}}$ 下特征值的不变性、状态转移矩阵 $\boldsymbol{\Phi}(t)=\mathrm{e}^{\mathbf{A}t}$ 与时变 $\boldsymbol{\Phi}(t,t_0)$，以及约旦标准型化简——这些是本章反复使用的工具。本章要回答上一章\emph{回避}的问题：状态空间描述比传递函数“看得更深”，深在哪里？答曰：深在它能分辨出传递函数\emph{看不见}的内部结构（不可控 / 不可观的模态）。可控性与可观性，正是这种“内部可视性”的精确刻画。

\begin{table}[H]\centering\small
\caption{本章阅读方式建议}
\begin{tabular}{lll}
\toprule
阅读方式 & 时间 & 适合谁 \\
\midrule
精读（含全部推导与例题） & 4--6 小时 & 首次系统学习、要做综合 / 估计推导 \\
速读（跳过冗长矩阵计算） & 1.5 小时 & 学过线性系统、想对齐本书记号与结构观 \\
速查（只看小结的符号表 + 判据速查表） & 15 分钟 & 设计反馈 / 观测器时回来查判据 \\
\bottomrule
\end{tabular}
\end{table}

\begin{note}
\textbf{本章记号与约定.}\;
状态 $\mathbf{x}\in\mathbb{R}^n$、输入 $\mathbf{u}\in\mathbb{R}^r$（多输入 $r>1$）、输出 $\mathbf{y}\in\mathbb{R}^m$，全书粗体（两本教材源用非粗体 / $\boldsymbol{u}$，融入时一律粗体化）。系统三元组记 $(\mathbf{A},\mathbf{B},\mathbf{C})$ 或 $\Sigma=(\mathbf{A},\mathbf{B},\mathbf{C})$；离散系统记 $(\mathbf{G},\mathbf{H},\mathbf{C})$（与导论 \cref{def:ctrl-controllable} 的离散记号 $(\mathbf{A},\mathbf{B})$ 同义，本章在离散专节用 $(\mathbf{G},\mathbf{H})$ 以贴合刘豹源、避免与连续混淆）。可控性矩阵记 $\mathcal{C}=[\mathbf{B}\ \mathbf{A}\mathbf{B}\ \cdots\ \mathbf{A}^{n-1}\mathbf{B}]$（导论 \cref{eq:ctrl-cmat} 同物，教材源记 $\mathbf{M}$）；可观性矩阵 $\mathcal{O}=[\mathbf{C};\mathbf{C}\mathbf{A};\cdots;\mathbf{C}\mathbf{A}^{n-1}]$（导论 \cref{eq:ctrl-omat}，教材源记 $\mathbf{N}$）。Gram 矩阵 $\mathbf{W}_c,\mathbf{W}_o$。状态转移矩阵 $\boldsymbol{\Phi}(t,t_0)$。\textbf{“可控”与“能控”、“可观”与“能观”在本书同义}（教材多用“能控 / 能观”，本书正文用“可控 / 可观”对齐导论，引用教材标题处保留“能控 / 能观”）。
\end{note}

% ════════════════════════════════════════════════════════════════
\section{可控性与可观性：问题、定义与历史}
\label{sec:co-def}

\paragraph{动机：状态空间带来的两个新问题。}
经典控制理论只关心输入到输出的传递关系，这关系由传递函数唯一确定；只要满足稳定性，输出就被认为是“可控”的，而输出本身可测、便也是“可观”的——经典理论根本无须区分这两件事\cite{liubao2006modern}。现代控制理论建立在\emph{状态空间}描述之上：状态方程刻画输入 $\mathbf{u}(t)$ 如何驱动\emph{内部}状态 $\mathbf{x}(t)$，输出方程刻画状态如何映成可测的 $\mathbf{y}(t)$。一旦把视线从“输入--输出”移到“输入--状态--输出”，立刻冒出两个经典理论从未提出的问题：

\begin{itemize}
  \item \textbf{可控性}：输入 $\mathbf{u}(t)$ 对状态 $\mathbf{x}(t)$ 的\emph{驱动能力}——能否用控制把状态从任意起点搬到任意终点？
  \item \textbf{可观性}：输出 $\mathbf{y}(t)$ 对状态 $\mathbf{x}(t)$ 的\emph{反映能力}——能否仅凭一段输出（与输入）反推出内部状态？
\end{itemize}

\noindent 这两个概念是\textbf{状态空间法的“原生概念”}，与“系统内部结构”一一对应；它们也正是后续一切设计的地基：状态反馈要求“能驱动”（可控），状态观测器要求“能反推”（可观），LQR / 卡尔曼滤波的解存在唯一性都直接建立在它们之上（见 \cref{sec:co-outlook}）。

\paragraph{历史：Kalman 1960 的两个发明。}
可控性与可观性由 R.\,E.\ Kalman 于 1960 年前后首次系统提出\cite{kalman1960general,kalman1963mathematical}，与他同期提出的卡尔曼滤波\cite{kalman1960new}、以及由此发展的最优控制 / 最优估计理论一脉相承——可控性是\emph{最优控制}的设计基础，可观性是\emph{最优估计}的设计基础，而二者经“对偶原理”（\cref{sec:co-duality}）精确相连。可以说，正是这两个概念把“状态空间”从一种记法升格为一套\emph{有判别力}的理论。

\begin{definition}[完全状态可控（连续定常）]\label{def:co-controllable}
对线性连续定常系统 $\dot{\mathbf{x}}=\mathbf{A}\mathbf{x}+\mathbf{B}\mathbf{u}$，若存在一个分段连续的输入 $\mathbf{u}(t)$，能在\emph{有限}时间区间 $[t_0,t_f]$ 内把系统从某初态 $\mathbf{x}(t_0)$ 转移到指定的\emph{任一}终态 $\mathbf{x}(t_f)$，则称该状态\textbf{可控}；若状态空间中\emph{所有}状态都可控，则称系统\textbf{状态完全可控}，简称\textbf{可控}\cite{liubao2006modern}。
\end{definition}

\noindent{}（这是导论 \cref{def:ctrl-controllable} 的展开版：导论给“存在输入把任意初态搬到任意终态”，本章额外强调“某状态可控”与“系统完全可控”之分——前者是单点性质，后者要求可控状态\emph{充满}整个状态空间。）

\begin{insight}
\textbf{可控 $\neq$ 稳定，是两件正交的事。}\;
新手最易混淆。“可控”问的是\emph{能否用输入改变状态}（一种“可操纵性”），“稳定”问的是\emph{放手不管时状态是否自行收敛}（一种“自治行为”）。一个不稳定系统可以完全可控（正因可控，才有救——可用状态反馈把它镇定，见 \cref{ch:linear-synthesis}）；一个稳定系统也可以不可控（某些模态虽自行衰减，却任凭你怎么加输入都纹丝不动）。本章只谈“能不能操纵 / 反推”，稳定性留给 \cref{ch:lyapunov-basics}。
\end{insight}

\paragraph{两条便利约定（定常系统专属）。}
在定常系统中，为简化讨论，可不失一般性地（\cref{def:co-controllable}）取以下两种等价提法之一\cite{liubao2006modern}：

\begin{enumerate}
  \item 设 $t_0=0$、终态指定为\emph{零}：判断能否把任意 $\mathbf{x}(0)$ 在有限时间内驱到 $\mathbf{x}(t_f)=\mathbf{0}$。这正是“可控”的提法。
  \item 设\emph{初态}为零：判断能否把 $\mathbf{x}(t_0)=\mathbf{0}$ 驱到任意终态 $\mathbf{x}(t_f)$。这叫\textbf{可达（reachable）}。
\end{enumerate}

\begin{proposition}[定常系统：可控 $\iff$ 可达]\label{prop:co-reach}
对线性连续定常系统，可控性与可达性等价：可控系统必可达，可达系统必可控\cite{liubao2006modern}。
\end{proposition}

\noindent 因此本章对定常系统不再区分二者，统称“可控”。（注意这一等价\emph{依赖} $\mathrm{e}^{\mathbf{A}t}$ 可逆这一定常专属事实；离散系统中 $\mathbf{G}$ 可奇异，可控与可达可能分裂——见 \cref{sec:co-discrete}。还需强调：可控性定义中控制作用\emph{理论上无约束}，我们只问“能否到达”，不计较状态轨迹的形状，也不计较 $\mathbf{u}$ 的大小。）

\paragraph{可观性的对称定义。}
可观性只关心输出反映状态的能力，与控制作用无直接关系；既然可由模型算出已知输入引起的那部分输出并扣除，分析可观性时不妨令 $\mathbf{u}=\mathbf{0}$，只从齐次状态方程与输出方程出发\cite{liubao2006modern}。

\begin{definition}[完全可观（连续定常）]\label{def:co-observable}
对系统 $\dot{\mathbf{x}}=\mathbf{A}\mathbf{x},\ \mathbf{y}=\mathbf{C}\mathbf{x}$（$\mathbf{x}(t_0)=\mathbf{x}_0$），若能根据\emph{有限}观测区间 $[t_0,t_f]$ 上的输出 $\mathbf{y}(t)$ \emph{唯一}地确定初态 $\mathbf{x}_0$，则称 $\mathbf{x}_0$ \textbf{可观}；若每个状态都可观，则称系统\textbf{状态完全可观}，简称\textbf{可观}\cite{liubao2006modern}。
\end{definition}

\begin{note}
\textbf{为何定义在“确定初态”上？}\;
一旦初态 $\mathbf{x}_0$ 确定，由状态转移方程 $\mathbf{x}(t)=\boldsymbol{\Phi}(t-t_0)\mathbf{x}_0+\int_{t_0}^t\boldsymbol{\Phi}(t-\tau)\mathbf{B}\mathbf{u}(\tau)\,\mathrm d\tau$，给定输入便可推得任意时刻的状态。故“能否反推全程状态”等价于“能否反推初态”——这就是为什么可观性只针对初态来定义\cite{liubao2006modern}。另：若输出维数 $m=n$ 且 $\mathbf{C}$ 非奇异，则 $\mathbf{x}=\mathbf{C}^{-1}\mathbf{y}$ 瞬间可解、根本不需观测时间；可观性问题之所以非平凡，正因为一般 $m<n$，须在多个时刻测量、拼出 $n$ 个独立方程。
\end{note}

% ════════════════════════════════════════════════════════════════
\section{结构直觉：约旦型与“孤立方块”}
\label{sec:co-jordan}

\paragraph{先建图像，再谈代数。}
秩判据（导论已给）是\emph{算法}，但它不解释“为什么”。要真正理解可控可观，最快的路径是把系统对角化 / 约旦化，让每个模态“独立站出来”——这时可控性退化成一个一眼可辨的结构问题。这一节是全章的物理图像之源。

\paragraph{对角型：一个 $b_i=0$ 就废一维。}
设系统已化为对角型 $\dot{\mathbf{x}}=\boldsymbol{\Lambda}\mathbf{x}+\mathbf{b}u$（单输入、特征值互异），逐分量写出即
\begin{equation}
  \dot{x}_i=\lambda_i x_i+b_i u,\qquad i=1,\dots,n.
  \label{eq:co-diag-scalar}
\end{equation}
这是 $n$ 个\emph{彼此解耦}的一阶标量方程——状态变量之间毫无联系，只能各自被输入 $u$ 直接驱动。于是结论一目了然\cite{liubao2006modern}：

\begin{center}\itshape
若 $b_i\neq 0$，第 $i$ 维受 $u$ 控制，可控；\\
若 $b_i=0$，第 $i$ 维成了齐次方程 $\dot{x}_i=\lambda_i x_i$，与 $u$ 无关，其解 $x_i(0)\mathrm{e}^{\lambda_i t}$ 无强迫分量、非零初值下永远无法在有限时间内归零——\textbf{不可控}。
\end{center}

\noindent 即：\textbf{对角型系统可控 $\iff$ $\mathbf{b}$（一般为 $\mathbf{T}^{-1}\mathbf{B}$）没有全零行}。可观性完全对称：把输出写成 $\mathbf{y}=\mathbf{C}\mathbf{x}$，对角型下 $\mathbf{y}(t)=\sum_i \mathbf{c}_{\cdot i}\,\mathrm{e}^{\lambda_i t}x_{i0}$（$\mathbf{c}_{\cdot i}$ 为 $\mathbf{C}$ 的第 $i$ 列）；若 $\mathbf{C}$ 的第 $i$ 列全为零，则输出根本不含 $x_i$ 的自由分量，$x_i$ 无从反推、\textbf{不可观}\cite{liubao2006modern}。即 \textbf{对角型可观 $\iff$ $\mathbf{C}$（一般为 $\mathbf{C}\mathbf{T}$）没有全零列}。

\paragraph{孤立方块：可控性的几何与结构图像。}
把 \cref{eq:co-diag-scalar} 画成模拟结构图，$b_i=0$ 的那一维就成了一个\emph{与输入毫无连线}的“孤立方块”：它对应特征值 $\lambda_i$ 的自然模式 $\mathrm{e}^{\lambda_i t}$，自顾自地自由运动，输入鞭长莫及。\cref{fig:co-isolated} 把“可控”与“不可控”两种结构并排画出——这张图是全章最值得记住的图像。

\begin{figure}[ht]
\centering
\begin{tikzpicture}[>=Stealth, font=\small,
  blk/.style={draw, rounded corners, minimum width=12mm, minimum height=8mm, align=center, fill=main!6, draw=main!60!black},
  iso/.style={draw, rounded corners, minimum width=12mm, minimum height=8mm, align=center, fill=third!10, draw=third!60!black, dashed},
  sum/.style={draw, circle, inner sep=1pt, minimum size=4mm},
  ln/.style={->, thick, gray!55!black}]
% ---- (a) controllable: B drives every block ----
\node at (-3.3, 2.55) {\textbf{(a) 可控}（无孤立方块）};
\node[blk] (a1) at (-4.6, 1.4) {$\dfrac{1}{s-\lambda_1}$};
\node[blk] (a2) at (-4.6, -0.1) {$\dfrac{1}{s-\lambda_2}$};
\coordinate (au) at (-6.4, 0.65);
\node[left] at (au) {$u$};
\draw[ln] (au) -- (-5.7,0.65) coordinate(asp);
\draw[ln] (asp) |- (a1.west) node[midway, above, font=\scriptsize]{$b_1\neq0$};
\draw[ln] (asp) |- (a2.west) node[midway, below, font=\scriptsize]{$b_2\neq0$};
\node[sum] (ay) at (-2.9,0.65) {$+$};
\draw[ln] (a1.east) -| (ay);
\draw[ln] (a2.east) -| (ay);
\draw[ln] (ay) -- ++(0.9,0) node[right]{$y$};
% ---- (b) uncontrollable: one isolated block ----
\node at (3.4, 2.55) {\textbf{(b) 不可控}（$x_1$ 孤立）};
\node[iso] (b1) at (2.1, 1.4) {$\dfrac{1}{s-\lambda_1}$};
\node[blk] (b2) at (2.1, -0.1) {$\dfrac{1}{s-\lambda_2}$};
\coordinate (bu) at (0.3, 0.0);
\node[left] at (bu) {$u$};
\draw[ln] (bu) -- (b2.west) node[midway, below, font=\scriptsize]{$b_2\neq0$};
\node[font=\scriptsize, third!60!black] at (1.05,1.0) {$b_1=0$（无连线）};
\fill (1.1,1.4) circle (0pt);
% isolated block has only an initial-condition input
\draw[ln, third!70!black] (1.05,2.05) node[above, font=\scriptsize]{$x_1(0)$} -- (b1.north);
\node[sum] (by) at (3.8,0.65) {$+$};
\draw[ln] (b1.east) -| (by);
\draw[ln] (b2.east) -| (by);
\draw[ln] (by) -- ++(0.9,0) node[right]{$y$};
\end{tikzpicture}
\caption{可控性的“孤立方块”图像。(a) 输入到每个模态都有连线（$b_i\neq0$），系统可控。(b) 第一个模态与输入\emph{无连线}（$b_1=0$），它成为一个仅由初值 $x_1(0)$ 激励、随后自由运动的“孤立方块”——其自然模式 $\mathrm{e}^{\lambda_1 t}$ 不可控。（仿刘豹\cite{liubao2006modern} 图 3.3 / 3.5 重绘）}
\label{fig:co-isolated}
\end{figure}

\paragraph{约旦块：可控性沿块“向上传染”。}
特征值有重根时系统化为约旦型。以二阶约旦块
\[
  \dot{\mathbf{x}}=\begin{pmatrix}\lambda_1 & 1\\ 0 & \lambda_1\end{pmatrix}\mathbf{x}+\begin{pmatrix}0\\ b_2\end{pmatrix}u
  \quad\Longleftrightarrow\quad
  \begin{cases}\dot{x}_1=\lambda_1 x_1+x_2,\\ \dot{x}_2=\lambda_1 x_2+b_2 u\end{cases}
\]
为例\cite{liubao2006modern}：$x_2$ 直接受 $u$ 控制（$b_2\neq0$）；$x_1$ 虽不直接连 $u$，却经约旦块的超对角“$1$”被 $x_2$ 驱动——\textbf{约旦块内前一个状态总被后一个状态拉着走}，故整块可控。反之，若把输入接在\emph{非末位}（如 $\mathbf{b}=(b_1,0)^\top$），则 $\dot{x}_2=\lambda_1 x_2$ 又成孤立齐次方程、不可控。由此得约旦型可控判据\cite{liubao2006modern}：

\begin{insight}
\textbf{约旦型可控 / 可观的“一行 / 一列”法则。}\;
\begin{itemize}
  \item \textbf{可控}：$\mathbf{T}^{-1}\mathbf{B}$ 中，与每个约旦块\emph{最后一行}相对应的那些行，元素不全为零（互异特征值部分则各行不全为零）。直觉：可控性从约旦块的“底”灌入、沿超对角线“向上”传遍全块。
  \item \textbf{可观}：$\mathbf{C}\mathbf{T}$ 中，与每个约旦块\emph{第一列}相对应的那些列，元素不全为零。直觉：可观性从约旦块的“顶”被输出“看到”、沿超对角线“向下”串联可见。
\end{itemize}
\textbf{陷阱（重根的致命例外）}：上述法则\emph{仅当每个重特征值只对应一个约旦块时成立}。若同一特征值有\emph{两个及以上}约旦块（即几何重数 $>1$），即便各块末行 / 首列非零也未必可控 / 可观——此时单输入必不可控，多输入要再查“那些末行元素构成的向量是否线性无关”\cite{liubao2006modern,zhang2017modern}。例如 $\mathbf{A}=\mathrm{diag}(1,1),\ \mathbf{b}=(1,1)^\top$：对角、$\mathbf{b}$ 无零行，却\emph{不可控}（两维模态被同一标量输入“同步”驱动，无法独立操纵）。
\end{insight}

% ════════════════════════════════════════════════════════════════
\section{秩判据：从无穷步到 \texorpdfstring{$n$}{n} 步}
\label{sec:co-rank}

约旦型直观但需先做相似变换。本节直接由 $(\mathbf{A},\mathbf{B})$ 给出秩判据，并把导论 \cref{eq:ctrl-cmat} 背后的推导走一遍——核心是看清 Cayley--Hamilton 定理如何把“无穷步可达”压成“$n$ 步可达”。

\subsection{Kalman 秩判据（连续，单输入）}
\label{sec:co-kalman-cont}

\begin{theorem}[Kalman 可控性秩判据]\label{thm:co-rank}
线性连续定常系统 $\dot{\mathbf{x}}=\mathbf{A}\mathbf{x}+\mathbf{b}u$（单输入）可控的充要条件是可控性矩阵
\begin{equation}
  \mathcal{C}=\begin{bmatrix}\mathbf{b} & \mathbf{A}\mathbf{b} & \mathbf{A}^2\mathbf{b} & \cdots & \mathbf{A}^{n-1}\mathbf{b}\end{bmatrix}\in\mathbb{R}^{n\times n}
  \label{eq:co-cmat}
\end{equation}
满秩，即 $\operatorname{rank}\mathcal{C}=n$\cite{liubao2006modern}。多输入情形 $\mathbf{b}\to\mathbf{B}\in\mathbb{R}^{n\times r}$，$\mathcal{C}\in\mathbb{R}^{n\times nr}$，判据仍为 $\operatorname{rank}\mathcal{C}=n$（见 \cref{sec:co-rank-multi}）。
\end{theorem}

\noindent{}（此即导论 \cref{eq:ctrl-cmat}；本章给完整推导。导论从离散迭代切入，这里从连续解的积分式切入，两条路殊途同归，互为印证。）

\begin{proof}[从连续解出发的推导]
单输入系统解为 $\mathbf{x}(t)=\boldsymbol{\Phi}(t-t_0)\mathbf{x}(t_0)+\int_{t_0}^t\boldsymbol{\Phi}(t-\tau)\mathbf{b}u(\tau)\,\mathrm d\tau$。令 $t=t_f$、$\mathbf{x}(t_f)=\mathbf{0}$，整理得可控状态须满足
\begin{equation}
  \mathbf{x}(t_0)=-\int_{t_0}^{t_f}\boldsymbol{\Phi}(t_0-\tau)\,\mathbf{b}\,u(\tau)\,\mathrm d\tau.
  \label{eq:co-control-relation}
\end{equation}
关键一步用 \textbf{Cayley--Hamilton 定理}：$\mathbf{A}$ 满足自身特征方程，故 $\mathbf{A}^k$（任意 $k$）可由 $\{\mathbf{I},\mathbf{A},\dots,\mathbf{A}^{n-1}\}$ 线性表出，于是
\begin{equation}
  \boldsymbol{\Phi}(t)=\mathrm{e}^{\mathbf{A}t}=\sum_{k=0}^{\infty}\frac{t^k}{k!}\mathbf{A}^k
  =\sum_{j=0}^{n-1}\beta_j(t)\,\mathbf{A}^j,
  \label{eq:co-CH-collapse}
\end{equation}
即\emph{无穷级数被压缩为 $n$ 项}，标量系数 $\beta_j(t)=\sum_{k=0}^{\infty}\alpha_{jk}\,t^k/k!$。代入 \cref{eq:co-control-relation}（单输入下 $u$ 为标量、积分为标量）：
\[
  \mathbf{x}(t_0)=-\sum_{j=0}^{n-1}\mathbf{A}^j\mathbf{b}\underbrace{\int_{t_0}^{t_f}\beta_j(t_0-\tau)u(\tau)\,\mathrm d\tau}_{=:\,\gamma_j}
  =-\begin{bmatrix}\mathbf{b}&\mathbf{A}\mathbf{b}&\cdots&\mathbf{A}^{n-1}\mathbf{b}\end{bmatrix}
  \begin{bmatrix}\gamma_0\\ \gamma_1\\ \vdots\\ \gamma_{n-1}\end{bmatrix}.
\]
要对\emph{任意} $\mathbf{x}(t_0)$ 都能解出标量组 $(\gamma_0,\dots,\gamma_{n-1})$，当且仅当系数矩阵 $\mathcal{C}=[\mathbf{b}\ \mathbf{A}\mathbf{b}\ \cdots\ \mathbf{A}^{n-1}\mathbf{b}]$ 可逆，即 $\operatorname{rank}\mathcal{C}=n$。判据得证。
\end{proof}

\begin{insight}
\textbf{为什么算到 $\mathbf{A}^{n-1}$ 就够？}\;（前置自测第 2 问的答案）正是 \cref{eq:co-CH-collapse}：Cayley--Hamilton 定理保证 $\mathbf{A}^n,\mathbf{A}^{n+1},\dots$ 都是 $\{\mathbf{I},\dots,\mathbf{A}^{n-1}\}$ 的线性组合，再往后算\emph{不会带来新的可达方向}。所以“$n$ 步内不可达，则任意步数皆不可达”——可达方向空间在 $n$ 步处“封顶”。这与导论从离散迭代得出的同一结论（\cref{eq:ctrl-reach} 附近）完全吻合，是同一数学事实的两种叙述。
\end{insight}

\begin{example}[可控标准型必可控]\label{ex:co-cstd}
考察 $\dot{\mathbf{x}}=\begin{pmatrix}0&1&0\\0&0&1\\-a_0&-a_1&-a_2\end{pmatrix}\mathbf{x}+\begin{pmatrix}0\\0\\1\end{pmatrix}u$。逐次计算
\[
  \mathbf{b}=\begin{pmatrix}0\\0\\1\end{pmatrix},\quad
  \mathbf{A}\mathbf{b}=\begin{pmatrix}0\\1\\-a_2\end{pmatrix},\quad
  \mathbf{A}^2\mathbf{b}=\begin{pmatrix}1\\-a_2\\a_2^2-a_1\end{pmatrix},
\]
故 $\mathcal{C}=\begin{pmatrix}0&0&1\\0&1&-a_2\\1&-a_2&a_2^2-a_1\end{pmatrix}$ 是反三角阵、副对角全为 $1$，无论 $a_i$ 取何值 $\operatorname{rank}\mathcal{C}=3$。\textbf{凡这种“友矩阵 + 末行单位输入”形式的系统恒可控}，故称\emph{可控标准型}（\cref{sec:co-canonical} 详述）\cite{liubao2006modern}。
\end{example}

\begin{example}[一个不可控系统：约旦法与秩法互校]\label{ex:co-uncontrol}
设 $\dot{\mathbf{x}}=\begin{pmatrix}-4&5\\1&0\end{pmatrix}\mathbf{x}+\begin{pmatrix}-5\\1\end{pmatrix}u$。\emph{约旦法}：特征值 $\lambda=-5,1$，取 $\mathbf{T}=[\mathbf{p}_1\ \mathbf{p}_2]=\begin{pmatrix}-5&1\\1&1\end{pmatrix}$，算得 $\mathbf{T}^{-1}\mathbf{b}=\begin{pmatrix}1\\0\end{pmatrix}$——第二行为零，对应模态 $\mathrm{e}^{t}$ 不可控。\emph{秩法}：$\mathcal{C}=[\mathbf{b}\ \mathbf{A}\mathbf{b}]=\begin{pmatrix}-5&25\\1&-5\end{pmatrix}$，$\det\mathcal{C}=0$，$\operatorname{rank}\mathcal{C}=1<2$，不可控。两法结论一致，且都指认是不稳定模态 $\mathrm{e}^{t}$ 失控——这种系统\emph{不可镇定}（见 \cref{sec:co-stabilizable}）\cite{liubao2006modern}。
\end{example}

\subsection{多输入与 Gramian 视角}
\label{sec:co-rank-multi}

\paragraph{多输入：行多于列时的实用算法。}
多输入情形 $\mathbf{B}\in\mathbb{R}^{n\times r}$，$\mathcal{C}=[\mathbf{B}\ \mathbf{A}\mathbf{B}\ \cdots\ \mathbf{A}^{n-1}\mathbf{B}]$ 是 $n\times nr$ 的“矮胖”矩阵，秩仍要求为 $n$。证明思路与单输入相同，只是 \cref{eq:co-control-relation} 里的标量 $\gamma_j$ 变成 $r$ 维向量 $\boldsymbol{\Gamma}_j$，方程组从“$n$ 个未知数 $n$ 个方程”变为“$nr$ 个未知数 $n$ 个方程”，由线性方程组理论，对任意右端有解 $\iff \operatorname{rank}\mathcal{C}=\operatorname{rank}[\mathcal{C}\,|\,\mathbf{x}(t_0)]\iff\mathcal{C}$ 满行秩\cite{liubao2006modern}。计算上，行比列少时常用 $\operatorname{rank}\mathcal{C}=\operatorname{rank}(\mathcal{C}\mathcal{C}^\top)$（化为 $n\times n$ 方阵判秩）。

\begin{note}
\textbf{多输入更“容易”可控。}\;
多输入系统往往无须用满 $\mathcal{C}$：常常前几个分块 $[\mathbf{B}\ \mathbf{A}\mathbf{B}]$ 就已满秩。直觉上，输入通道越多、能“撬动”的状态方向越多，可控条件越容易满足；极端地，若 $r=n$ 且 $\mathbf{B}$ 可逆，一步即可任意配置状态\cite{liubao2006modern}。这也是为什么机器人（多关节力矩输入）通常天然可控。
\end{note}

\paragraph{Gramian 判据（与导论的接口）。}
导论 \cref{thm:ctrl-gramian-rank} 已给出等价的 Gramian 判据：可控 $\iff$ 可控性 Gramian $\mathbf{W}_c\succ0$，其中 $\mathbf{W}_c$ 满足 Lyapunov 方程 $\mathbf{A}\mathbf{W}_c+\mathbf{W}_c\mathbf{A}^\top=-\mathbf{B}\mathbf{B}^\top$（\cref{eq:ctrl-gramian}），有限时域形式 $\mathbf{W}_c(t_0,t_1)=\int_{t_0}^{t_1}\boldsymbol{\Phi}(t_0,t)\mathbf{B}\mathbf{B}^\top\boldsymbol{\Phi}(t_0,t)^\top\,\mathrm dt$。本章不重复其证明（导论已给双向构造性证明，含最小能量输入 $\mathbf{u}(t)=\mathbf{B}^\top\boldsymbol{\Phi}^\top\mathbf{W}_c^{-1}(\cdots)$），但在 \cref{sec:co-timevarying} 把它推广到时变系统——届时会看到，\textbf{Gramian 判据才是“母判据”，秩判据是它在定常情形的特例}。

\subsection{PBH 判据与可镇定 / 可检测}
\label{sec:co-stabilizable}

\paragraph{PBH：逐模态检验。}
导论 \cref{eq:ctrl-pbh} 给出 Popov--Belevitch--Hautus（PBH）判据：可控 $\iff\operatorname{rank}[\lambda\mathbf{I}-\mathbf{A}\ \ \mathbf{B}]=n,\ \forall\lambda\in\mathbb{C}$（实际只需对 $\mathbf{A}$ 的特征值检验）。它与秩判据等价，但视角更“按模态”：在特征值 $\lambda_i$ 处，$[\lambda_i\mathbf{I}-\mathbf{A}]$ 降秩、其左零向量就是\emph{该模态对应的左特征向量}，PBH 要求这个左特征向量\emph{不}与 $\mathbf{B}$ 的所有列正交——即“该模态能被某个输入通道‘触及’”。这与 \cref{sec:co-jordan} 的孤立方块图像是同一回事的代数版。可观 PBH 对偶：$\operatorname{rank}\big[\begin{smallmatrix}\lambda\mathbf{I}-\mathbf{A}\\ \mathbf{C}\end{smallmatrix}\big]=n,\ \forall\lambda$。

\begin{insight}
\textbf{可镇定 / 可检测：只管“坏模态”。}\;
完全可控是“每个模态都能操纵”，但工程上往往\emph{不必}——稳定的模态放任自流即可，只需\emph{不稳定}的模态能被操纵。这就是\textbf{可镇定（stabilizable）}：PBH 仅对 $\operatorname{Re}\lambda\ge0$（连续）或 $|\lambda|\ge1$（离散）的 $\lambda$ 成立。对偶地，\textbf{可检测（detectable）}只要求不稳定模态可观。这两个\emph{松弛}概念是后续理论的真正前提：
\begin{itemize}
  \item LQR 解（CARE 唯一正定）只需 $(\mathbf{A},\mathbf{B})$ \emph{可镇定}（不必完全可控）——见 \cref{ch:lqr-lqg-riccati}；
  \item 卡尔曼滤波 / 观测器收敛只需 $(\mathbf{A},\mathbf{C})$ \emph{可检测}——见 \cref{ch:eskf}。
\end{itemize}
\cref{ex:co-uncontrol} 的系统失控的恰是\emph{不稳定}模态 $\mathrm{e}^t$，故\emph{不可镇定}——这是“可控”与“可镇定”分野的最小反例：去掉它，系统虽不可控但仍可镇定。
\end{insight}

% ════════════════════════════════════════════════════════════════
\section{可观性与对偶原理}
\label{sec:co-observe}

\subsection{可观性矩阵}
\label{sec:co-omat}

可观性的推导与可控性\emph{完全对称}。由 $\mathbf{y}(t)=\mathbf{C}\mathbf{x}(t)=\mathbf{C}\,\mathrm{e}^{\mathbf{A}(t-t_0)}\mathbf{x}_0$，再用 \cref{eq:co-CH-collapse} 把 $\mathrm{e}^{\mathbf{A}(t-t_0)}$ 收缩为 $n$ 项 $\sum_j\beta_j\mathbf{A}^j$，得
\begin{equation}
  \mathbf{y}(t)=\sum_{j=0}^{n-1}\beta_j(t-t_0)\,\mathbf{C}\mathbf{A}^j\,\mathbf{x}_0
  =\big(\beta_0\mathbf{I}\ \beta_1\mathbf{I}\ \cdots\ \beta_{n-1}\mathbf{I}\big)
  \underbrace{\begin{bmatrix}\mathbf{C}\\ \mathbf{C}\mathbf{A}\\ \vdots\\ \mathbf{C}\mathbf{A}^{n-1}\end{bmatrix}}_{=:\,\mathcal{O}}\mathbf{x}_0.
  \label{eq:co-omat-derive}
\end{equation}
要从一段输出唯一定出 $\mathbf{x}_0$，当且仅当\textbf{可观性矩阵}
\begin{equation}
  \mathcal{O}=\begin{bmatrix}\mathbf{C}\\ \mathbf{C}\mathbf{A}\\ \vdots\\ \mathbf{C}\mathbf{A}^{n-1}\end{bmatrix}\in\mathbb{R}^{mn\times n}
  \quad\text{满秩},\qquad \text{可观}\iff\operatorname{rank}\mathcal{O}=n,
  \label{eq:co-omat}
\end{equation}
此即导论 \cref{eq:ctrl-omat}\cite{liubao2006modern}。不可观子空间 $\mathcal{N}=\bigcap_{k=0}^{n-1}\ker(\mathbf{C}\mathbf{A}^k)=\ker\mathcal{O}$；可观 $\iff\mathcal{N}=\{\mathbf{0}\}$。

\subsection{对偶原理}
\label{sec:co-duality}

\paragraph{对偶：把一套结论用两次。}
可观性矩阵 \cref{eq:co-omat} 与可控性矩阵 \cref{eq:co-cmat} 长得像“互相转置”，这并非巧合，而是 Kalman 的\textbf{对偶原理}\cite{kalman1960general,liubao2006modern}。把它讲透，可观性的全部判据都能从可控性\emph{白嫖}。

\begin{definition}[对偶系统]\label{def:co-dual}
称 $\Sigma_2=(\mathbf{A}_2,\mathbf{B}_2,\mathbf{C}_2)$ 是 $\Sigma_1=(\mathbf{A}_1,\mathbf{B}_1,\mathbf{C}_1)$ 的\textbf{对偶系统}，若
\begin{equation}
  \mathbf{A}_2=\mathbf{A}_1^\top,\qquad \mathbf{B}_2=\mathbf{C}_1^\top,\qquad \mathbf{C}_2=\mathbf{B}_1^\top.
  \label{eq:co-dual-def}
\end{equation}
几何上：对偶系统把输入端与输出端互换、信号传递方向反转、引出点与综合点互换、各矩阵转置（\cref{fig:co-duality}）\cite{liubao2006modern}。
\end{definition}

\begin{figure}[ht]
\centering
\begin{tikzpicture}[>=Stealth, font=\small,
  blk/.style={draw, rounded corners, minimum width=14mm, minimum height=9mm, align=center, fill=main!6, draw=main!60!black},
  ln/.style={->, thick, gray!55!black}]
% Sigma_1
\node at (-3.2,1.7) {$\Sigma_1=(\mathbf{A}_1,\mathbf{B}_1,\mathbf{C}_1)$};
\node[blk] (s1B) at (-4.5,0) {$\mathbf{B}_1$};
\node[blk] (s1A) at (-3.2,0) {$(s\mathbf{I}-\mathbf{A}_1)^{-1}$};
\node[blk] (s1C) at (-1.7,0) {$\mathbf{C}_1$};
\draw[ln] (-5.6,0) node[left]{$\mathbf{u}_1$} -- (s1B);
\draw[ln] (s1B) -- (s1A);
\draw[ln] (s1A) -- (s1C);
\draw[ln] (s1C) -- ++(1.0,0) node[right]{$\mathbf{y}_1$};
% Sigma_2 (arrows reversed, transposed)
\node at (3.2,1.7) {$\Sigma_2=(\mathbf{A}_1^\top,\mathbf{C}_1^\top,\mathbf{B}_1^\top)$};
\node[blk] (s2C) at (4.5,0) {$\mathbf{B}_1^\top$};
\node[blk] (s2A) at (3.2,0) {$(s\mathbf{I}-\mathbf{A}_1^\top)^{-1}$};
\node[blk] (s2B) at (1.7,0) {$\mathbf{C}_1^\top$};
\draw[ln] (5.6,0) node[right]{$\mathbf{u}_2$} -- (s2C);
\draw[ln] (s2C) -- (s2A);
\draw[ln] (s2A) -- (s2B);
\draw[ln] (s2B) -- ++(-1.0,0) node[left]{$\mathbf{y}_2$};
\end{tikzpicture}
\caption{对偶系统：输入 / 输出端互换、信号流反向、各矩阵转置。$\Sigma_1$ 的可控性 = $\Sigma_2$ 的可观性。（仿刘豹\cite{liubao2006modern} 图 3.11 重绘）}
\label{fig:co-duality}
\end{figure}

\begin{theorem}[对偶原理]\label{thm:co-duality}
若 $\Sigma_1,\Sigma_2$ 互为对偶，则 $\Sigma_1$ 的可控性等价于 $\Sigma_2$ 的可观性，$\Sigma_1$ 的可观性等价于 $\Sigma_2$ 的可控性\cite{liubao2006modern}。等价地：$(\mathbf{A},\mathbf{B})$ 可控 $\iff(\mathbf{A}^\top,\mathbf{B}^\top)$ 可观；$(\mathbf{A},\mathbf{C})$ 可观 $\iff(\mathbf{A}^\top,\mathbf{C}^\top)$ 可控。
\end{theorem}

\noindent{}（此即导论 \cref{thm:ctrl-duality}；本章给一行证明。）
\begin{proof}
$\Sigma_2$ 的可控性矩阵 $\mathcal{C}_2=[\mathbf{B}_2\ \mathbf{A}_2\mathbf{B}_2\ \cdots\ \mathbf{A}_2^{n-1}\mathbf{B}_2]$。代入 \cref{eq:co-dual-def}（$\mathbf{A}_2=\mathbf{A}_1^\top,\ \mathbf{B}_2=\mathbf{C}_1^\top$）：
\[
  \mathcal{C}_2=\big[\mathbf{C}_1^\top\ \ \mathbf{A}_1^\top\mathbf{C}_1^\top\ \ \cdots\ \ (\mathbf{A}_1^\top)^{n-1}\mathbf{C}_1^\top\big]
  =\mathcal{O}_1^\top,
\]
其中 $\mathcal{O}_1=[\mathbf{C}_1;\mathbf{C}_1\mathbf{A}_1;\cdots;\mathbf{C}_1\mathbf{A}_1^{n-1}]$ 是 $\Sigma_1$ 的可观性矩阵。转置不改变秩，故 $\operatorname{rank}\mathcal{C}_2=\operatorname{rank}\mathcal{O}_1$，即 $\Sigma_2$ 可控 $\iff\Sigma_1$ 可观。另一半同理（用 $\mathcal{O}_2^\top=\mathcal{C}_1$）。
\end{proof}

\begin{insight}
\textbf{对偶是“省力工具”，但物理意义相反。}\;
把 $\mathbf{A}\to\mathbf{A}^\top,\ \mathbf{B}\to\mathbf{C}^\top$ 一替换，可控性矩阵 \cref{eq:co-cmat} 立刻变成可观性矩阵 \cref{eq:co-omat} 的转置，可控 Gramian 变可观 Gramian，PBH 上下叠放。于是\emph{可控性的每条结论都自动有一条可观性版本}，本章后文的标准型、结构分解都将“做一遍可控、对偶得可观”。\textbf{边界（务必牢记）}：对偶是\emph{数学}上的镜像，两侧物理意义\emph{相反}——一个问“能否驱动”、一个问“能否观测”。这条对偶在 \cref{ch:eskf} 再度现身：LQR 增益与卡尔曼增益由\emph{对偶}的两个 Riccati 给出，“最优控制”与“最优估计”由此打通（导论 \cref{sec:ctrl-lqe-dual} 已点明，C 章 \cref{ch:lqr-lqg-riccati} 给完整证明）。
\end{insight}

\begin{note}
\textbf{时变对偶的细微差别。}\;
时变系统的对偶定义要在 $\mathbf{A}_2(t)=-\mathbf{A}_1^\top(t)$ 处\emph{多一个负号}（$\mathbf{B}_2=\mathbf{C}_1^\top,\ \mathbf{C}_2=\mathbf{B}_1^\top$）。由此可推出两对偶系统的状态转移矩阵互为\emph{转置逆}：$\boldsymbol{\Phi}_2^\top(t_0,t)=\boldsymbol{\Phi}_1(t,t_0)$，进而 $\Sigma_2$ 的可控 Gram 矩阵恰等于 $\Sigma_1$ 的可观 Gram 矩阵（$\mathbf{W}_{c2}=\mathbf{W}_{o1}$），对偶原理依然成立\cite{liubao2006modern}。定常情形 $\mathbf{A}^\top$ 与 $-\mathbf{A}^\top$ 的特征值仅差符号、不影响秩，故 \cref{eq:co-dual-def} 用 $\mathbf{A}^\top$ 即可。
\end{note}

% ════════════════════════════════════════════════════════════════
\section{时变系统的 Gram 判据}
\label{sec:co-timevarying}

\paragraph{为什么需要它：机器人沿轨迹线性化是 LTV。}
机器人动力学 $\mathbf{M}(\mathbf{q})\ddot{\mathbf{q}}+\mathbf{C}(\mathbf{q},\dot{\mathbf{q}})\dot{\mathbf{q}}+\mathbf{g}(\mathbf{q})=\boldsymbol{\tau}$ 沿一条标称轨迹线性化，得到的是\emph{时变}线性系统（LTV）$\dot{\mathbf{x}}=\mathbf{A}(t)\mathbf{x}+\mathbf{B}(t)\mathbf{u}$——系统矩阵随时间变化。这时 $(\mathbf{A},\mathbf{B})$ 是常数矩阵的秩判据失效，必须改用 Gram 矩阵\cite{liubao2006modern,zhang2017modern}。本节是 \cref{sec:co-rank-multi} Gramian 视角的时变推广，也是 iLQR / 轨迹优化（\cref{ch:optimal-control-basics}）背后“沿轨迹可控”的理论基础。

\begin{theorem}[时变可控性 Gram 判据]\label{thm:co-tv-gram}
时变系统 $\dot{\mathbf{x}}=\mathbf{A}(t)\mathbf{x}+\mathbf{B}(t)\mathbf{u}$ 在 $[t_0,t_f]$ 上状态完全可控的充要条件是 Gram 矩阵
\begin{equation}
  \mathbf{W}_c(t_0,t_f)=\int_{t_0}^{t_f}\boldsymbol{\Phi}(t_0,t)\,\mathbf{B}(t)\mathbf{B}^\top(t)\,\boldsymbol{\Phi}^\top(t_0,t)\,\mathrm dt
  \label{eq:co-tv-gramc}
\end{equation}
非奇异（$\boldsymbol{\Phi}(t,t_0)$ 为系统的状态转移矩阵）\cite{liubao2006modern}。
\end{theorem}

\begin{proof}[要点]
\emph{充分性（构造）}：设 $\mathbf{W}_c$ 非奇异，取控制 $\mathbf{u}(t)=-\mathbf{B}^\top(t)\boldsymbol{\Phi}^\top(t_0,t)\mathbf{W}_c^{-1}(t_0,t_f)\mathbf{x}(t_0)$，代入解的积分式直接验证 $\mathbf{x}(t_f)=\mathbf{0}$——故任意初态可驱到原点。\emph{必要性（反证）}：若 $\mathbf{W}_c$ 奇异，存在 $\mathbf{x}(t_0)\neq\mathbf{0}$ 使 $\mathbf{x}^\top(t_0)\mathbf{W}_c\mathbf{x}(t_0)=\int_{t_0}^{t_f}\|\mathbf{B}^\top(t)\boldsymbol{\Phi}^\top(t_0,t)\mathbf{x}(t_0)\|^2\mathrm dt=0$，由被积函数连续推出 $\mathbf{B}^\top(t)\boldsymbol{\Phi}^\top(t_0,t)\mathbf{x}(t_0)\equiv\mathbf{0}$；再用可控状态须满足的关系式 \cref{eq:co-control-relation}（时变版）导出 $\|\mathbf{x}(t_0)\|^2=0$，与 $\mathbf{x}(t_0)\neq\mathbf{0}$ 矛盾\cite{liubao2006modern}。
\end{proof}

\begin{insight}
\textbf{Gram 判据是“母判据”。}\;
为何说时变 Gram 判据统一了定常秩判据？因为定常时 $\boldsymbol{\Phi}(t_0,t)=\mathrm{e}^{\mathbf{A}(t_0-t)}$，再用 \cref{eq:co-CH-collapse} 把 $\mathrm{e}^{\mathbf{A}(t_0-t)}\mathbf{B}$ 展成 $[\mathbf{B}\ \mathbf{A}\mathbf{B}\ \cdots\ \mathbf{A}^{n-1}\mathbf{B}]\,[\beta_0;\dots;\beta_{n-1}]$，于是“$\mathbf{W}_c$ 非奇异 $\iff$ $\mathbf{B}^\top\boldsymbol{\Phi}^\top$ 行向量线性无关 $\iff\operatorname{rank}\mathcal{C}=n$”\cite{liubao2006modern}。一句话：\textbf{时变与定常判据形异实同、一脉相承，Gram 矩阵是 $(\mathbf{A},\mathbf{B})$ 秩判据的一般形式}。
\end{insight}

\begin{note}
\textbf{实用充分判据（免算 STM）。}\;
\cref{thm:co-tv-gram} 需先求出状态转移矩阵 $\boldsymbol{\Phi}$，但 LTV 的 $\boldsymbol{\Phi}$ 往往无闭式。一个只用 $\mathbf{A}(t),\mathbf{B}(t)$ 的\emph{充分}判据是：令 $\mathbf{B}_1(t)=\mathbf{B}(t)$、$\mathbf{B}_i(t)=-\mathbf{A}(t)\mathbf{B}_{i-1}(t)+\dot{\mathbf{B}}_{i-1}(t)$（$i=2,\dots,n$），$\mathbf{Q}_c(t)=[\mathbf{B}_1\ \mathbf{B}_2\ \cdots\ \mathbf{B}_n]$；若存在某 $t_f$ 使 $\operatorname{rank}\mathbf{Q}_c(t_f)=n$，则系统在 $[0,t_f]$ 可控\cite{liubao2006modern}。可观侧对偶：$\mathbf{C}_1=\mathbf{C}$、$\mathbf{C}_i=\mathbf{C}_{i-1}\mathbf{A}+\dot{\mathbf{C}}_{i-1}$。注意这\emph{只是充分}条件——不满足不代表不可控。
\end{note}

% ════════════════════════════════════════════════════════════════
\section{可控标准型与可观标准型}
\label{sec:co-canonical}

\paragraph{动机：为设计选一副“好坐标”。}
状态变量的选取不唯一，状态空间表达式也不唯一；既然非奇异变换不改变可控性 / 可观性（\cref{ch:statespace-basics} 已证特征值不变，这里再加“结构不变”），就该挑一副\emph{对设计最方便}的坐标。事实证明\cite{liubao2006modern}：状态反馈用\textbf{可控标准型}最省事（极点配置可一眼写出，见 \cref{ch:linear-synthesis}），观测器设计 / 系统辨识用\textbf{可观标准型}最方便。只有完全可控（可观）的系统才能化为可控（可观）标准型。

\paragraph{可控标准型（友矩阵形）。}
单输入单输出系统若可控，存在非奇异变换将其化为如下\textbf{可控标准型}（友矩阵 / companion form）\cite{zhang2017modern}：
\begin{equation}
  \mathbf{A}_c=\begin{bmatrix}
    0&1&0&\cdots&0\\
    0&0&1&\cdots&0\\
    \vdots&\vdots&\vdots&\ddots&\vdots\\
    0&0&0&\cdots&1\\
    -a_n&-a_{n-1}&-a_{n-2}&\cdots&-a_1
  \end{bmatrix},\quad
  \mathbf{b}_c=\begin{bmatrix}0\\0\\\vdots\\0\\1\end{bmatrix},
  \label{eq:co-cstd}
\end{equation}
其中 $s^n+a_1 s^{n-1}+\cdots+a_n$ 是 $\mathbf{A}$ 的特征多项式。\cref{ex:co-cstd} 已验证此型恒可控；其妙处在于\textbf{特征多项式系数直接显式排在末行}——状态反馈 $\mathbf{u}=-\mathbf{k}\bar{\mathbf{x}}$ 时，$\mathbf{A}_c-\mathbf{b}_c\mathbf{k}$ 仍是友矩阵、末行变为 $-(a_i+k_i)$，于是“配极点”就是“配末行”，闭环特征多项式系数可逐项凑出（Ackermann 公式的来由，\cref{ch:linear-synthesis} 详述）。

\paragraph{可观标准型（对偶得来）。}
由对偶原理，可观标准型几乎是可控标准型的“转置版”。SISO 可观标准型一种常见形式为
\begin{equation}
  \mathbf{A}_o=\begin{bmatrix}
    0&0&\cdots&0&-a_n\\
    1&0&\cdots&0&-a_{n-1}\\
    0&1&\cdots&0&-a_{n-2}\\
    \vdots&\vdots&\ddots&\vdots&\vdots\\
    0&0&\cdots&1&-a_1
  \end{bmatrix},\quad
  \mathbf{c}_o=\begin{bmatrix}0&0&\cdots&0&1\end{bmatrix},
  \label{eq:co-ostd}
\end{equation}
特征多项式系数排在\emph{末列}\cite{liubao2006modern}。它对观测器设计的便利与可控标准型对状态反馈的便利完全对偶。

\subsection{显式变换阵 \texorpdfstring{$\mathbf{T}_{c1}$}{Tc1} 与 \texorpdfstring{$\boldsymbol{\beta}_i$}{beta} 输出系数（能控标准 I 型）}
\label{sec:co-canon-transform}

\paragraph{补上“怎么变过去”这一环。}
上面给出的 $(\mathbf{A}_c,\mathbf{b}_c)$ 与 $(\mathbf{A}_o,\mathbf{c}_o)$ 是\emph{目标形貌}，但还缺最关键的一步：给定一个一般的、绝不长成友矩阵样子的 $(\mathbf{A},\mathbf{b},\mathbf{c})$，那个把它\emph{搬过去}的非奇异变换阵到底是什么、输出阵 $\bar{\mathbf{c}}$ 的元素又由谁决定？这一小节把变换阵 $\mathbf{T}_{c1}$ 一笔一画地构造出来，并给出 $\bar{\mathbf{c}}=(\beta_0,\dots,\beta_{n-1})$ 的显式公式——它正是 \cref{sec:syn-ackermann} 极点配置、\cref{sec:co-minreal} 最小实现真正动手时所依赖的“坐标手术刀”。

\begin{theorem}[能控标准 I 型的变换阵与输出系数]\label{thm:co-canon-transform}
设单输入系统 $\dot{\mathbf{x}}=\mathbf{A}\mathbf{x}+\mathbf{b}u,\ y=\mathbf{c}\mathbf{x}$ 完全可控，其特征多项式为 $|s\mathbf{I}-\mathbf{A}|=s^n+a_1 s^{n-1}+\cdots+a_{n-1}s+a_n$。则线性变换 $\mathbf{x}=\mathbf{T}_{c1}\bar{\mathbf{x}}$，
\begin{equation}
  \mathbf{T}_{c1}=\big[\mathbf{e}_1\ \ \mathbf{e}_2\ \cdots\ \mathbf{e}_n\big],\qquad
  \mathbf{e}_{n-k}=\sum_{j=0}^{k}a_j\,\mathbf{A}^{k-j}\mathbf{b}\quad(a_0:=1,\ k=0,\dots,n-1),
  \label{eq:co-tc1}
\end{equation}
把系统化为能控标准 I 型 \cref{eq:co-cstd}：$\bar{\mathbf{A}}=\mathbf{T}_{c1}^{-1}\mathbf{A}\mathbf{T}_{c1}=\mathbf{A}_c$，$\bar{\mathbf{b}}=\mathbf{T}_{c1}^{-1}\mathbf{b}=(0,\dots,0,1)^\top=\mathbf{b}_c$，而输出阵
\begin{equation}
  \bar{\mathbf{c}}=\mathbf{c}\,\mathbf{T}_{c1}=(\beta_0,\ \beta_1,\ \dots,\ \beta_{n-1}),\qquad
  \beta_k=\mathbf{c}\,\mathbf{e}_{k+1}=\mathbf{c}\sum_{j=0}^{\,n-1-k}a_j\,\mathbf{A}^{\,n-1-k-j}\mathbf{b}.
  \label{eq:co-beta1}
\end{equation}
此 $\beta_i$ 恰是变换后传递函数 $W(s)=\bar{\mathbf{c}}(s\mathbf{I}-\bar{\mathbf{A}})^{-1}\bar{\mathbf{b}}=\dfrac{\beta_{n-1}s^{n-1}+\cdots+\beta_1 s+\beta_0}{s^n+a_1 s^{n-1}+\cdots+a_n}$ 的分子系数\cite{liubao2006modern}。
\end{theorem}

\begin{proof}[构造法（取自刘豹 §3.7.1，含线性无关性）]
因系统可控，$\{\mathbf{b},\mathbf{A}\mathbf{b},\dots,\mathbf{A}^{n-1}\mathbf{b}\}$ 线性无关；\cref{eq:co-tc1} 把它们做\emph{上三角}线性组合（$\mathbf{e}_n=\mathbf{b}$，$\mathbf{e}_{n-1}=\mathbf{A}\mathbf{b}+a_1\mathbf{b}$，$\dots$，$\mathbf{e}_1=\mathbf{A}^{n-1}\mathbf{b}+a_1\mathbf{A}^{n-2}\mathbf{b}+\cdots+a_{n-1}\mathbf{b}$），系数矩阵对角线全为 $1$ 故可逆，于是 $\mathbf{e}_1,\dots,\mathbf{e}_n$ 仍线性无关，$\mathbf{T}_{c1}$ 非奇异。关键是验证 $\mathbf{A}\mathbf{T}_{c1}=\mathbf{T}_{c1}\mathbf{A}_c$，即逐列算 $\mathbf{A}\mathbf{e}_j$。对 $k=0,\dots,n-2$，
\[
  \mathbf{A}\mathbf{e}_{n-k}=\sum_{j=0}^{k}a_j\mathbf{A}^{k+1-j}\mathbf{b}
  =\Big(\sum_{j=0}^{k+1}a_j\mathbf{A}^{k+1-j}\mathbf{b}\Big)-a_{k+1}\mathbf{b}
  =\mathbf{e}_{n-k-1}-a_{k+1}\mathbf{e}_n,
\]
即每作用一次 $\mathbf{A}$ 就“前移一列”、并漏下一个 $-a_{k+1}\mathbf{e}_n$。唯独最高列 $\mathbf{e}_1$ 越界，须用 Cayley--Hamilton（\cref{eq:co-CH-collapse} 的同源事实）$\mathbf{A}^{n}=-a_1\mathbf{A}^{n-1}-\cdots-a_n\mathbf{I}$ 收口：
\[
  \mathbf{A}\mathbf{e}_1=\big(\mathbf{A}^{n}+a_1\mathbf{A}^{n-1}+\cdots+a_{n-1}\mathbf{A}\big)\mathbf{b}=-a_n\mathbf{b}=-a_n\mathbf{e}_n.
\]
把这 $n$ 个等式按列排开，$\mathbf{A}\mathbf{e}_1,\dots,\mathbf{A}\mathbf{e}_n$ 写成 $\mathbf{T}_{c1}$ 各列的组合，其系数阵正是友矩阵 $\mathbf{A}_c$（末列来自 $\mathbf{A}\mathbf{e}_1=-a_n\mathbf{e}_n$、其余来自“前移”），故 $\bar{\mathbf{A}}=\mathbf{A}_c$。又 $\mathbf{b}=\mathbf{e}_n=\mathbf{T}_{c1}(0,\dots,0,1)^\top$ 给出 $\bar{\mathbf{b}}=(0,\dots,0,1)^\top$；$\bar{\mathbf{c}}=\mathbf{c}\mathbf{T}_{c1}=(\mathbf{c}\mathbf{e}_1,\dots,\mathbf{c}\mathbf{e}_n)$ 即 \cref{eq:co-beta1}。
\end{proof}

\paragraph{把 \texorpdfstring{$\mathbf{T}_{c1}$}{Tc1} 写成“可控性矩阵 $\times$ Toeplitz 阵”。}
\cref{eq:co-tc1} 的上三角组合可一次性写成矩阵乘积——这是手算与编程时真正在用的形式：
\begin{equation}
  \mathbf{T}_{c1}=\big[\mathbf{A}^{n-1}\mathbf{b}\ \ \mathbf{A}^{n-2}\mathbf{b}\ \cdots\ \mathbf{A}\mathbf{b}\ \ \mathbf{b}\big]
  \begin{bmatrix}
    1 & & & & \\
    a_1 & 1 & & & \\
    a_2 & a_1 & 1 & & \\
    \vdots & & \ddots & \ddots & \\
    a_{n-1} & \cdots & a_2 & a_1 & 1
  \end{bmatrix}.
  \label{eq:co-tc1-toeplitz}
\end{equation}
右乘的下三角 Toeplitz 阵（沿对角线元素恒定、首列为 $(1,a_1,\dots,a_{n-1})^\top$）只由特征多项式系数搭成。这与 \cref{sec:co-rank-multi} 的可控性矩阵 \cref{eq:co-cmat}（只是列序倒置）相乘即得 $\mathbf{T}_{c1}$。\textbf{它正是 Bass--Gura 算法（\cref{algo:syn-bassgura}）第 3 步那个 $\mathbf{Q}=[\mathbf{b}\ \mathbf{A}\mathbf{b}\ \cdots]\cdot\mathbf{W}$ 的“运算版”}——极点配置时无须显式化型，但其代数骨架就是这把 $\mathbf{T}_{c1}$。

\begin{example}[一般 $(\mathbf{A},\mathbf{b},\mathbf{c})\to$ 能控标准 I 型]\label{ex:co-canon-transform}
取 $\mathbf{A}=\big(\begin{smallmatrix}1&1\\4&-2\end{smallmatrix}\big),\ \mathbf{b}=\big(\begin{smallmatrix}1\\0\end{smallmatrix}\big),\ \mathbf{c}=(0,1)$。特征多项式 $|s\mathbf{I}-\mathbf{A}|=s^2+s-6$，故 $a_1=1,\ a_2=-6$。先算 $\mathbf{A}\mathbf{b}=(1,4)^\top$，可控性矩阵 $\mathcal{C}=[\mathbf{b}\ \mathbf{A}\mathbf{b}]=\big(\begin{smallmatrix}1&1\\0&4\end{smallmatrix}\big)$ 满秩。按 \cref{eq:co-tc1-toeplitz}（$n=2$，Toeplitz 阵 $\big(\begin{smallmatrix}1&0\\a_1&1\end{smallmatrix}\big)=\big(\begin{smallmatrix}1&0\\1&1\end{smallmatrix}\big)$）：
\[
  \mathbf{T}_{c1}=[\mathbf{A}\mathbf{b}\ \ \mathbf{b}]\begin{pmatrix}1&0\\1&1\end{pmatrix}
  =\begin{pmatrix}1&1\\4&0\end{pmatrix}\begin{pmatrix}1&0\\1&1\end{pmatrix}=\begin{pmatrix}2&1\\4&0\end{pmatrix}.
\]
于是 $\bar{\mathbf{A}}=\mathbf{T}_{c1}^{-1}\mathbf{A}\mathbf{T}_{c1}=\big(\begin{smallmatrix}0&1\\6&-1\end{smallmatrix}\big)$（末行 $(-a_2,-a_1)=(6,-1)$，与 \cref{eq:co-cstd} 吻合），$\bar{\mathbf{b}}=(0,1)^\top$。输出系数由 \cref{eq:co-beta1}：$\beta_0=\mathbf{c}\mathbf{e}_1=\mathbf{c}(\mathbf{A}\mathbf{b}+a_1\mathbf{b})=(0,1)(2,4)^\top=4$，$\beta_1=\mathbf{c}\mathbf{e}_2=\mathbf{c}\mathbf{b}=0$，故 $\bar{\mathbf{c}}=\mathbf{c}\mathbf{T}_{c1}=(4,0)$。回读传递函数 $W(s)=\dfrac{\beta_1 s+\beta_0}{s^2+a_1 s+a_2}=\dfrac{4}{s^2+s-6}$，与直接由 $\mathbf{c}(s\mathbf{I}-\mathbf{A})^{-1}\mathbf{b}$ 算得的完全一致——可控性保证了这条“化型即读系数”的捷径成立。
\end{example}

\subsection{能控标准 II 型：变换阵就是可控性矩阵本身}
\label{sec:co-canon-II}

\paragraph{第二种约定，更省事的一型。}
同一可控系统还有第二种标准型，区别只在“选哪 $n$ 个无关列、按何顺序排”。\textbf{能控标准 II 型}干脆取变换阵为可控性矩阵\emph{原样}\cite{liubao2006modern}：
\begin{equation}
  \mathbf{x}=\mathbf{T}_{c2}\bar{\mathbf{x}},\qquad
  \mathbf{T}_{c2}=\big[\mathbf{b}\ \ \mathbf{A}\mathbf{b}\ \cdots\ \mathbf{A}^{n-1}\mathbf{b}\big]=\mathcal{C}.
  \label{eq:co-tc2}
\end{equation}
此变换无须任何 Toeplitz 修饰，最便于手算。代价是 $\bar{\mathbf{A}}$ 把特征多项式系数排到\emph{末列}而非末行：
\begin{equation}
  \bar{\mathbf{A}}=\mathbf{T}_{c2}^{-1}\mathbf{A}\mathbf{T}_{c2}=\begin{bmatrix}
    0&0&\cdots&0&-a_n\\
    1&0&\cdots&0&-a_{n-1}\\
    0&1&\cdots&0&-a_{n-2}\\
    \vdots&\vdots&\ddots&\vdots&\vdots\\
    0&0&\cdots&1&-a_1
  \end{bmatrix},\quad
  \bar{\mathbf{b}}=\begin{bmatrix}1\\0\\\vdots\\0\end{bmatrix},\quad
  \bar{\mathbf{c}}=\mathbf{c}\mathbf{T}_{c2}=(\beta_0,\dots,\beta_{n-1}),\ \ \beta_i=\mathbf{c}\mathbf{A}^i\mathbf{b}.
  \label{eq:co-tc2-form}
\end{equation}
推导一行即明：$\mathbf{A}\mathbf{T}_{c2}=[\mathbf{A}\mathbf{b}\ \cdots\ \mathbf{A}^{n}\mathbf{b}]$，再用 Cayley--Hamilton 把越界的 $\mathbf{A}^{n}\mathbf{b}=-a_1\mathbf{A}^{n-1}\mathbf{b}-\cdots-a_n\mathbf{b}$ 回代，便得末列含 $-a_i$ 的友矩阵；$\bar{\mathbf{b}}$ 因 $\mathbf{b}$ 是 $\mathbf{T}_{c2}$ 首列而取 $(1,0,\dots,0)^\top$；$\bar{\mathbf{c}}$ 的元素 $\beta_i=\mathbf{c}\mathbf{A}^i\mathbf{b}$ 正是系统的\emph{马尔可夫参数}（脉冲响应矩阵的各阶矩），物理意义比 I 型的 $\beta_i$ 更直接\cite{liubao2006modern}。沿用 \cref{ex:co-canon-transform} 的系统，$\mathbf{T}_{c2}=\mathcal{C}=\big(\begin{smallmatrix}1&1\\0&4\end{smallmatrix}\big)$，$\bar{\mathbf{A}}=\big(\begin{smallmatrix}0&6\\1&-1\end{smallmatrix}\big)$，$\bar{\mathbf{b}}=(1,0)^\top$，$\bar{\mathbf{c}}=(\mathbf{c}\mathbf{b},\mathbf{c}\mathbf{A}\mathbf{b})=(0,4)$——与 I 型同源、同一 $W(s)$，矩阵却全然不同。

\paragraph{“一行生成”变换阵：直接造 \texorpdfstring{$\mathbf{T}_{c1}^{-1}$}{Tc1 inverse}。}
张嗣瀛给出一条与 I 型等价、却绕开 Toeplitz 的捷径\cite{zhang2017modern}：把目标写成 $\bar{\mathbf{x}}=\mathbf{P}\mathbf{x}$（即 $\mathbf{P}=\mathbf{T}_{c1}^{-1}$），只需取可控性矩阵\emph{逆阵的末行} $\mathbf{p}_1$，再逐次右乘 $\mathbf{A}$ 堆叠：
\begin{equation}
  \mathbf{p}_1=\big[\,0\ \cdots\ 0\ \ 1\,\big]\,\mathcal{C}^{-1},\qquad
  \mathbf{P}=\begin{bmatrix}\mathbf{p}_1\\ \mathbf{p}_1\mathbf{A}\\ \vdots\\ \mathbf{p}_1\mathbf{A}^{n-1}\end{bmatrix}=\mathbf{T}_{c1}^{-1}.
  \label{eq:co-zhang-P}
\end{equation}
其来由是：欲使 $\mathbf{P}\mathbf{A}\mathbf{P}^{-1}=\mathbf{A}_c$（友矩阵）且 $\mathbf{P}\mathbf{b}=(0,\dots,0,1)^\top$，把 $\mathbf{P}=[\mathbf{p}_1;\dots;\mathbf{p}_n]^{}$ 代入比对行，立得 $\mathbf{p}_{i+1}=\mathbf{p}_i\mathbf{A}$、且 $\mathbf{p}_1[\mathbf{b}\ \mathbf{A}\mathbf{b}\ \cdots\ \mathbf{A}^{n-1}\mathbf{b}]=[0\ \cdots\ 0\ 1]$，解出的 $\mathbf{p}_1$ 即 \cref{eq:co-zhang-P}\cite{zhang2017modern}。对 \cref{ex:co-canon-transform} 的系统，$\mathcal{C}^{-1}=\big(\begin{smallmatrix}1&-1/4\\0&1/4\end{smallmatrix}\big)$，末行 $\mathbf{p}_1=(0,1/4)$，$\mathbf{p}_1\mathbf{A}=(1,-1/2)$，故 $\mathbf{P}=\big(\begin{smallmatrix}0&1/4\\1&-1/2\end{smallmatrix}\big)=\mathbf{T}_{c1}^{-1}$，与 \cref{ex:co-canon-transform} 的 $\mathbf{T}_{c1}=\big(\begin{smallmatrix}2&1\\4&0\end{smallmatrix}\big)$ 互逆，殊途同归。

\begin{insight}
\textbf{友矩阵为何“天生可控”——频域一眼看穿。}\;
\cref{ex:co-cstd} 用秩判据验证了能控标准型可控，这里给出张嗣瀛的频域佐证\cite{zhang2017modern}：对友矩阵 $\mathbf{A}_c$ 与 $\mathbf{b}_c=(0,\dots,0,1)^\top$，$(s\mathbf{I}-\mathbf{A}_c)^{-1}\mathbf{b}_c=\dfrac{1}{\Delta(s)}\,(1,\ s,\ s^2,\ \dots,\ s^{n-1})^\top$（$\Delta(s)$ 为特征多项式）。即各状态恰是 $x_k(s)=s^{k-1}u(s)/\Delta(s)$——后一个状态是前一个的\emph{导数}，故称\textbf{相变量（phase variable）}。取 $\mathbf{c}=(1,0,\dots,0)$ 便有 $W(s)=1/\Delta(s)$，分子分母\emph{绝无}公因子可消，由 \cref{thm:co-cancel} 即知必可控且可观。这把“友矩阵可控”从代数事实升格为“相变量链条节节相连”的物理图像。
\end{insight}

\begin{insight}
\textbf{I 型 vs.\ II 型：两套约定，别混用。}\;
同一可控系统有\emph{两副}能控标准型，差别仅在变换阵：I 型用 Toeplitz 修饰的 $\mathbf{T}_{c1}$（\cref{eq:co-tc1-toeplitz}）、系数落在 $\mathbf{A}_c$ \emph{末行}，便于和传递函数分母\emph{逐项}对照（\cref{sec:syn-ackermann} 极点配置取此型）；II 型用可控性矩阵原样 $\mathbf{T}_{c2}=\mathcal{C}$（\cref{eq:co-tc2}）、系数落在\emph{末列}、$\bar{\mathbf{c}}$ 即马尔可夫参数，构造最省力（手算判秩后顺带得到）。两者经一次置换相联，对 SISO 都唯一。\textbf{记号陷阱}：刘豹源以 $a_0,\dots,a_{n-1}$ 编号（常数项 $a_0$ 落角），本书随张嗣瀛取 $a_1,\dots,a_n$（常数项 $a_n$ 落角），二者只是系数下标\emph{倒序}，公式形不变、对号入座即可。
\end{insight}

\subsection{能观标准型的两型（由对偶得来）}
\label{sec:co-canon-obs}

可观侧完全对偶，亦分 I、II 两型，各与能控 I、II 型互为转置\cite{liubao2006modern}。设系统可观（$\operatorname{rank}\mathcal{O}=n$），\textbf{能观标准 I 型}的变换阵之逆\emph{就是可观性矩阵}：
\begin{equation}
  \mathbf{T}_{o1}^{-1}=\mathcal{O}=\begin{bmatrix}\mathbf{c}\\ \mathbf{c}\mathbf{A}\\ \vdots\\ \mathbf{c}\mathbf{A}^{n-1}\end{bmatrix}
  \ \Longrightarrow\
  \tilde{\mathbf{A}}=\mathbf{T}_{o1}^{-1}\mathbf{A}\mathbf{T}_{o1}=\mathbf{A}_c,\quad
  \tilde{\mathbf{b}}=\mathbf{T}_{o1}^{-1}\mathbf{b}=(\beta_0,\dots,\beta_{n-1})^\top,\quad
  \tilde{\mathbf{c}}=\mathbf{c}\mathbf{T}_{o1}=(1,0,\dots,0).
  \label{eq:co-to1}
\end{equation}
它恰是\emph{对偶系统}$(\mathbf{A}^\top,\mathbf{c}^\top,\mathbf{b}^\top)$ 能控 II 型的转置（\cref{thm:co-duality}）——可控的“变换阵 $=\mathcal{C}$”对偶成可观的“逆变换阵 $=\mathcal{O}$”，干净利落。\textbf{能观标准 II 型}则是能控 I 型的转置，其形貌正是本节开头展示的 \cref{eq:co-ostd}（系数排末列、$\tilde{\mathbf{c}}=(0,\dots,0,1)$、$\tilde{\mathbf{b}}=(\beta_0,\dots,\beta_{n-1})^\top$），变换阵 $\mathbf{T}_{o2}^{-1}$ 为可观性矩阵左乘一个上三角 Toeplitz 阵\cite{liubao2006modern}。可观 I 型的 $\tilde{\mathbf{c}}=(1,0,\dots,0)$ 把“第一个状态直接露在输出里”，是观测器与系统辨识最顺手的坐标。

\begin{note}
\textbf{为何能观型把 $\beta_i$ 挪到 $\tilde{\mathbf{b}}$ 上？}\;
对偶把输入端与输出端对调（\cref{def:co-dual}）：能控型里“输入怎样灌进各状态”由 $\bar{\mathbf{b}}$ 的标准列承担、特征信息落在 $\bar{\mathbf{c}}=(\beta_i)$；转置后角色互换，能观型里 $\tilde{\mathbf{c}}$ 取标准行、特征信息落到 $\tilde{\mathbf{b}}=(\beta_i)^\top$。\cref{eq:co-beta1} 那套 $\beta_i$ 公式原样搬用，无须重算\cite{liubao2006modern}。这再次印证 \cref{sec:co-duality} 的总纲：可控性的每条结论都“免费”附赠一条可观性版本。
\end{note}

\begin{pitfall}[多变量陷阱：可观标准型 $\neq$ 可控标准型的简单转置]
对 SISO 系统，可控标准型与可观标准型确实互为转置。但\textbf{对多输入多输出系统，二者绝非简单转置关系}\cite{liubao2006modern}！多变量传递函数阵 $\mathbf{W}(s)=\dfrac{\boldsymbol{\beta}_{n-1}s^{n-1}+\cdots+\boldsymbol{\beta}_0}{s^n+a_{n-1}s^{n-1}+\cdots+a_0}$（$\boldsymbol{\beta}_i$ 为 $m\times r$ 常数阵）的可控标准型实现维数是 $nr$（$\mathbf{A}_c$ 用 $r\times r$ 的 $\mathbf{0}_r,\mathbf{I}_r$ 分块、末行系数 $-a_i\mathbf{I}_r$，$\mathbf{C}_c=[\boldsymbol{\beta}_0\ \cdots\ \boldsymbol{\beta}_{n-1}]$），而可观标准型实现维数是 $nm$（分块用 $m\times m$，$\mathbf{B}_0=[\boldsymbol{\beta}_0;\cdots;\boldsymbol{\beta}_{n-1}]$，$\mathbf{C}_0=[\mathbf{0}_m\ \cdots\ \mathbf{0}_m\ \mathbf{I}_m]$）——\emph{维数都未必相等}（$nr$ vs $nm$），何谈转置。务必当心，不要把 SISO 的便利错推到 MIMO。
\end{pitfall}

% ════════════════════════════════════════════════════════════════
\section{Kalman 结构分解}
\label{sec:co-decomp}

\paragraph{把系统“拆成四块”。}
若系统不完全可控、不完全可观，所有可控状态构成\emph{可控子空间}、所有可观状态构成\emph{可观子空间}，但一般形式下这些子空间并未显式分离。本节用非奇异（坐标）变换把状态空间按可控性与可观性\textbf{结构分解}，揭示其本质结构——这正是最小实现问题（\cref{sec:co-minreal}）的理论依据，也与状态反馈、系统镇定密切相关\cite{liubao2006modern}。

\subsection{按可控性分解}
\label{sec:co-decomp-c}

\begin{theorem}[可控性分解]\label{thm:co-decomp-c}
若 $\operatorname{rank}\mathcal{C}=n_1<n$，则存在非奇异变换 $\mathbf{x}=\mathbf{R}_c\hat{\mathbf{x}}$，把 $(\mathbf{A},\mathbf{B},\mathbf{C})$ 变为
\begin{equation}
  \hat{\mathbf{A}}=\mathbf{R}_c^{-1}\mathbf{A}\mathbf{R}_c=\begin{bmatrix}\hat{\mathbf{A}}_{11}&\hat{\mathbf{A}}_{12}\\ \mathbf{0}&\hat{\mathbf{A}}_{22}\end{bmatrix},\quad
  \hat{\mathbf{B}}=\mathbf{R}_c^{-1}\mathbf{B}=\begin{bmatrix}\hat{\mathbf{B}}_1\\ \mathbf{0}\end{bmatrix},\quad
  \hat{\mathbf{C}}=\mathbf{C}\mathbf{R}_c=\begin{bmatrix}\hat{\mathbf{C}}_1&\hat{\mathbf{C}}_2\end{bmatrix},
  \label{eq:co-decomp-c}
\end{equation}
其中前 $n_1$ 维子系统 $\dot{\hat{\mathbf{x}}}_1=\hat{\mathbf{A}}_{11}\hat{\mathbf{x}}_1+\hat{\mathbf{B}}_1\mathbf{u}+\hat{\mathbf{A}}_{12}\hat{\mathbf{x}}_2$ 可控，后 $(n-n_1)$ 维子系统 $\dot{\hat{\mathbf{x}}}_2=\hat{\mathbf{A}}_{22}\hat{\mathbf{x}}_2$ 不可控\cite{liubao2006modern}。
\end{theorem}

\noindent\textbf{变换阵构造}：$\mathbf{R}_c=[\mathbf{R}_1\ \cdots\ \mathbf{R}_{n_1}\ \cdots\ \mathbf{R}_n]$ 的前 $n_1$ 列取\emph{可控性矩阵 $\mathcal{C}$ 中 $n_1$ 个线性无关的列}，其余 $(n-n_1)$ 列在保证 $\mathbf{R}_c$ 非奇异的前提下\emph{任取}\cite{liubao2006modern}。结构上，$\hat{\mathbf{B}}$ 的下半为零意味着\emph{输入 $\mathbf{u}$ 触不到 $\hat{\mathbf{x}}_2$}，$\hat{\mathbf{x}}_2$ 只作无控自由运动；$\hat{\mathbf{A}}$ 的下三角块为零意味着\emph{不可控部分不“回灌”可控部分}（但可控部分受不可控部分经 $\hat{\mathbf{A}}_{12}$ 影响）。丢掉 $\hat{\mathbf{x}}_2$ 即得一个低维可控系统。

\begin{example}[按可控性分解]\label{ex:co-decomp-c}
设 $\dot{\mathbf{x}}=\begin{pmatrix}0&0&-1\\1&0&-3\\0&1&-3\end{pmatrix}\mathbf{x}+\begin{pmatrix}1\\1\\0\end{pmatrix}u,\ \mathbf{y}=(0,1,-2)\mathbf{x}$。可控性矩阵 $\mathcal{C}=[\mathbf{b}\ \mathbf{A}\mathbf{b}\ \mathbf{A}^2\mathbf{b}]=\begin{pmatrix}1&0&-1\\1&1&-3\\0&1&-2\end{pmatrix}$，$\operatorname{rank}\mathcal{C}=2<3$，不完全可控。取 $\mathbf{R}_c=[\mathbf{b}\ \ \mathbf{A}\mathbf{b}\ \ \mathbf{R}_3]=\begin{pmatrix}1&0&0\\1&1&0\\0&1&1\end{pmatrix}$（$\mathbf{R}_3$ 任取保非奇异），变换后
\[
  \dot{\hat{\mathbf{x}}}=\begin{pmatrix}0&-1&-1\\1&-2&-2\\0&0&-1\end{pmatrix}\hat{\mathbf{x}}+\begin{pmatrix}1\\0\\0\end{pmatrix}u,\quad
  \mathbf{y}=(1,-1,-2)\hat{\mathbf{x}},
\]
清楚分出二维可控子系统 $\big(\begin{smallmatrix}0&-1\\1&-2\end{smallmatrix}\big),(1,0)^\top$（恰为可控标准型）与一维不可控子系统 $\dot{\hat{x}}_3=-\hat{x}_3$\cite{liubao2006modern}。
\end{example}

\subsection{按可观性分解与四块标准分解}
\label{sec:co-decomp-co}

按可观性分解由对偶得到：若 $\operatorname{rank}\mathcal{O}=n_1<n$，存在 $\mathbf{x}=\mathbf{R}_o\tilde{\mathbf{x}}$ 使 $\tilde{\mathbf{A}}$ 呈\emph{下}三角分块、$\tilde{\mathbf{C}}=[\tilde{\mathbf{C}}_1\ \mathbf{0}]$，分出可观子系统与不可观子系统（$\mathbf{R}_o^{-1}$ 的前 $n_1$ \emph{行}取 $\mathcal{O}$ 的 $n_1$ 个线性无关行）\cite{liubao2006modern}。把两次分解\emph{合起来}，得到本章的结构高峰：

\begin{theorem}[Kalman 四块标准分解]\label{thm:co-kalman-decomp}
任意 LTI 系统经一次非奇异变换 $\mathbf{x}=\mathbf{R}\bar{\mathbf{x}}$ 可分解为四部分——可控可观 $\bar{\mathbf{x}}_{co}$、可控不可观 $\bar{\mathbf{x}}_{c\bar{o}}$、不可控可观 $\bar{\mathbf{x}}_{\bar{c}o}$、不可控不可观 $\bar{\mathbf{x}}_{\bar{c}\bar{o}}$\cite{liubao2006modern}：
\begin{equation}
  \bar{\mathbf{A}}=\mathbf{R}^{-1}\mathbf{A}\mathbf{R}=\begin{bmatrix}
    \mathbf{A}_{11}&\mathbf{0}&\mathbf{A}_{13}&\mathbf{0}\\
    \mathbf{A}_{21}&\mathbf{A}_{22}&\mathbf{A}_{23}&\mathbf{A}_{24}\\
    \mathbf{0}&\mathbf{0}&\mathbf{A}_{33}&\mathbf{0}\\
    \mathbf{0}&\mathbf{0}&\mathbf{A}_{43}&\mathbf{A}_{44}
  \end{bmatrix},\
  \bar{\mathbf{B}}=\begin{bmatrix}\mathbf{B}_1\\ \mathbf{B}_2\\ \mathbf{0}\\ \mathbf{0}\end{bmatrix},\
  \bar{\mathbf{C}}=\begin{bmatrix}\mathbf{C}_1&\mathbf{0}&\mathbf{C}_3&\mathbf{0}\end{bmatrix},
  \label{eq:co-kalman-4block}
\end{equation}
其中 $(\mathbf{A}_{11},\mathbf{B}_1,\mathbf{C}_1)$ 是可控可观子系统。
\end{theorem}

\noindent{}（这里四块下标用 $c/\bar c$ 表可控 / 不可控、$o/\bar o$ 表可观 / 不可观——这套上 / 下划线记号在教材 OCR 稿中部分丢失，本书据结构补全。）\cref{fig:co-kalman} 把四块的信息流画出：\textbf{从输入 $\mathbf{u}$ 到输出 $\mathbf{y}$ 只有一条单向通道 $\mathbf{u}\to\mathbf{B}_1\to(\mathbf{A}_{11})\to\mathbf{C}_1\to\mathbf{y}$，途经可控且可观的那一块}。

\begin{figure}[ht]
\centering
\begin{tikzpicture}[>=Stealth, font=\small,
  box/.style={draw, rounded corners, minimum width=20mm, minimum height=11mm, align=center},
  co/.style={box, fill=main!10, draw=main!65!black, very thick},
  cobar/.style={box, fill=second!8, draw=second!60!black},
  cbaro/.style={box, fill=second!8, draw=second!60!black},
  dead/.style={box, fill=gray!12, draw=gray!55!black, dashed},
  ln/.style={->, thick, gray!55!black},
  lnm/.style={->, very thick, main!65!black}]
\node[co]    (co)    at (0,0)    {$\mathbf{x}_{co}$\\\scriptsize 可控·可观};
\node[cobar] (cobar) at (0,-2.0) {$\mathbf{x}_{c\bar o}$\\\scriptsize 可控·不可观};
\node[cbaro] (cbaro) at (4.2,0)  {$\mathbf{x}_{\bar c o}$\\\scriptsize 不可控·可观};
\node[dead]  (dead)  at (4.2,-2.0){$\mathbf{x}_{\bar c\bar o}$\\\scriptsize 不可控·不可观};
\coordinate (u) at (-3.6,0);
\node[left] at (u) {$\mathbf{u}$};
\draw[lnm] (u) -- node[above, font=\scriptsize]{$\mathbf{B}_1$} (co.west);
\draw[ln]  (-3.6,-2.0) node[left]{} -- node[below, font=\scriptsize]{$\mathbf{B}_2$} (cobar.west);
\coordinate (y) at (7.8,0);
\node[right] at (y) {$\mathbf{y}$};
\draw[lnm] (co.east) -- node[above, font=\scriptsize]{$\mathbf{C}_1$} (y);
\draw[ln]  (cbaro.east) -- node[above, font=\scriptsize]{$\mathbf{C}_3$} (7.8,0.0|-cbaro);
\draw[ln, dashed] (cbaro.north) ++(0,0) -- ++(0,0.0);
% internal couplings (lower-triangular flavor)
\draw[ln] (co.south) -- node[left, font=\scriptsize]{$\mathbf{A}_{21}$} (cobar.north);
\draw[ln] (cbaro.west) -- node[above, font=\scriptsize]{$\mathbf{A}_{13}$} (co.east);
\draw[ln] (cbaro.south) -- node[right, font=\scriptsize]{$\mathbf{A}_{43}$} (dead.north);
\node[main!65!black, font=\scriptsize, align=center] at (2.1,-3.2) {\textbf{唯一输入--输出通道}：$\mathbf{u}\to\mathbf{B}_1\to(\mathbf{A}_{11})\to\mathbf{C}_1\to\mathbf{y}$};
\end{tikzpicture}
\caption{Kalman 四块结构分解的信息流。传递函数只“看得见”可控且可观的子块 $(\mathbf{A}_{11},\mathbf{B}_1,\mathbf{C}_1)$；另外三块（灰 / 蓝）对输入--输出关系\emph{隐形}。（仿刘豹\cite{liubao2006modern} 图 3.14 重绘）}
\label{fig:co-kalman}
\end{figure}

\begin{theorem}[传递函数只反映可控可观子系统]\label{thm:co-tf-co}
对分解后的系统，传递函数阵
\begin{equation}
  \mathbf{W}(s)=\mathbf{C}(s\mathbf{I}-\mathbf{A})^{-1}\mathbf{B}=\mathbf{C}_1(s\mathbf{I}-\mathbf{A}_{11})^{-1}\mathbf{B}_1,
  \label{eq:co-tf-co}
\end{equation}
即\textbf{只取决于可控且可观的子块} $(\mathbf{A}_{11},\mathbf{B}_1,\mathbf{C}_1)$\cite{liubao2006modern}。
\end{theorem}

\begin{insight}
\textbf{传递函数是“有损”的内部描述。}\;
\cref{thm:co-tf-co} 是本章最深刻的论断之一：在系统里\emph{添加或删去}不可控 / 不可观的子系统，\textbf{丝毫不改变传递函数}。换言之，传递函数对内部结构是“看不全”的——它只反映那条唯一的输入--输出通道。这立刻引出两个推论：(1) 同一传递函数对应\emph{无穷多个}内部结构各异的状态空间实现（多出来的部分都是不可控 / 不可观的“暗物质”）；(2) 其中\emph{维数最小}的那个实现最有价值——它正好不含任何冗余的不可控 / 不可观部分。这就是\textbf{最小实现}（\cref{sec:co-minreal}）。\textbf{工程警示}：那些对传递函数“隐形”的不可控 / 不可观模态并非无害——若其中含不稳定模态，系统会在你看不见的地方发散（经典理论的零极点对消正是这种隐患的源头，见 \cref{sec:co-cancel}）。
\end{insight}

% ════════════════════════════════════════════════════════════════
\section{传递函数的实现与最小实现}
\label{sec:co-minreal}

\paragraph{实现问题：从外部描述反求内部。}
给定传递函数阵 $\mathbf{W}(s)$，若状态空间表达式 $\Sigma=(\mathbf{A},\mathbf{B},\mathbf{C},\mathbf{D})$ 满足 $\mathbf{C}(s\mathbf{I}-\mathbf{A})^{-1}\mathbf{B}+\mathbf{D}=\mathbf{W}(s)$，则称 $\Sigma$ 为 $\mathbf{W}(s)$ 的一个\textbf{实现}\cite{liubao2006modern}。物理可实现要求 $\mathbf{W}(s)$ 各元为\emph{真有理分式}（分子次数 $\le$ 分母次数）；若某元分子次数\emph{等于}分母次数，则实现含直接传输项 $\mathbf{D}$，且
\begin{equation}
  \mathbf{D}=\lim_{s\to\infty}\mathbf{W}(s)
  \label{eq:co-D-limit}
\end{equation}
（高频极限给出 $\mathbf{D}$；先扣除 $\mathbf{D}$ 使 $\mathbf{W}(s)-\mathbf{D}$ 严格真，再对其求 $(\mathbf{A},\mathbf{B},\mathbf{C})$ 实现）\cite{liubao2006modern}。

\begin{note}
\textbf{OCR 勘误.}\;
教材原 OCR 稿此处误作 $\mathbf{D}=\lim_{s\to 0}\mathbf{W}(s)$；正确应为 $s\to\infty$（直接传输项是\emph{高频}极限，$s\to0$ 给出的是稳态增益、与 $\mathbf{D}$ 无关）。本书据 \cref{eq:co-D-limit} 取正确版（与刘豹源例 3-18 实际算出的 $\mathbf{D}$ 自洽）。
\end{note}

\begin{definition}[最小实现]\label{def:co-minreal}
$\mathbf{W}(s)$ 的实现 $(\mathbf{A},\mathbf{B},\mathbf{C})$ 称为\textbf{最小实现}，若不存在维数更低的实现\cite{liubao2006modern}。
\end{definition}

\begin{theorem}[最小实现 $\iff$ 可控且可观]\label{thm:co-minreal}
$\mathbf{W}(s)$ 的实现 $\Sigma=(\mathbf{A},\mathbf{B},\mathbf{C})$ 为最小实现的\emph{充要条件}是 $\Sigma$ \textbf{既可控又可观}\cite{liubao2006modern}。
\end{theorem}

\noindent 这是 \cref{thm:co-tf-co} 的直接推论：既然传递函数只反映可控可观子块，任何不可控 / 不可观分量都是可删的冗余，删到只剩可控可观子块即达最小。\textbf{求最小实现的标准流程}\cite{liubao2006modern}：(1) 先写出任一实现（最方便取可控标准型或可观标准型）；(2) 对它做结构分解、抽出可控且可观子块 $(\hat{\mathbf{A}}_{11},\hat{\mathbf{B}}_1,\hat{\mathbf{C}}_1)$，即为最小实现。

\begin{example}[一个传递函数阵的最小实现]\label{ex:co-minreal}
求 $\mathbf{W}(s)=\Big(\dfrac{1}{(s+1)(s+2)}\ \ \dfrac{1}{(s+2)(s+3)}\Big)$ 的最小实现（单输出 $m=1$、双输入 $r=2$）。通分写成降幂标准格式 $\mathbf{W}(s)=\dfrac{(1,1)s+(3,1)}{s^3+6s^2+11s+6}$，读出 $a_0=6,a_1=11,a_2=6$、$\boldsymbol{\beta}_0=(3,1),\boldsymbol{\beta}_1=(1,1),\boldsymbol{\beta}_2=(0,0)$。先取可观标准型实现
\[
  \mathbf{A}_0=\begin{pmatrix}0&0&-6\\1&0&-11\\0&1&-6\end{pmatrix},\
  \mathbf{B}_0=\begin{pmatrix}3&1\\1&1\\0&0\end{pmatrix},\
  \mathbf{C}_0=(0,0,1).
\]
检验其可控性：$\mathcal{C}=[\mathbf{B}_0\ \mathbf{A}_0\mathbf{B}_0\ \mathbf{A}_0^2\mathbf{B}_0]$，算得 $\operatorname{rank}\mathcal{C}=3=n$。故 $(\mathbf{A}_0,\mathbf{B}_0,\mathbf{C}_0)$ 既可控又可观，\emph{本身即最小实现}（$n=3$）\cite{liubao2006modern}。
\end{example}

\begin{insight}
\textbf{最小实现唯一到“代数等价”。}\;
同一 $\mathbf{W}(s)$ 的实现不唯一，最小实现也不唯一（\cref{ex:co-minreal} 若改从可控标准型出发，会得到另一个数值不同的三维实现）。但可证：\textbf{任意两个最小实现必经某状态变换 $\mathbf{x}=\mathbf{P}\hat{\mathbf{x}}$ 相联}，即 $\tilde{\mathbf{A}}_m=\mathbf{P}^{-1}\mathbf{A}_m\mathbf{P},\ \tilde{\mathbf{B}}_m=\mathbf{P}^{-1}\mathbf{B}_m,\ \tilde{\mathbf{C}}_m=\mathbf{C}_m\mathbf{P}$——\emph{代数等价}\cite{liubao2006modern}。所以\textbf{最小实现的“维数”是唯一的}（等于 $\mathbf{W}(s)$ 的 McMillan 度），具体矩阵则随坐标而异。这个“维数唯一、坐标自由”正是后续状态估计 / 辨识中“可辨识性”概念的根（子空间辨识恢复的就是某个最小实现，见 \cref{ch:robust-hinf}）。
\end{insight}

% ════════════════════════════════════════════════════════════════
\section{零极点对消：可控可观的经典面孔（SISO）}
\label{sec:co-cancel}

\paragraph{把“现代”焊回“经典”。}
既然可控且可观 $\iff$ 最小实现，能否直接\emph{从传递函数表面}看出可控可观？对单输入单输出系统，答案极其干净\cite{liubao2006modern}：

\begin{theorem}[SISO：可控可观 $\iff$ 无零极点对消]\label{thm:co-cancel}
单输入单输出系统 $\Sigma=(\mathbf{A},\mathbf{b},\mathbf{c})$ 既可控又可观的\emph{充要条件}是其传递函数 $W(s)=\mathbf{c}(s\mathbf{I}-\mathbf{A})^{-1}\mathbf{b}$ 的\textbf{分子分母之间没有零极点对消}\cite{liubao2006modern}。
\end{theorem}

\begin{proof}[要点]
\emph{必要性}（逆否）：若 $\Sigma$ 不是最小实现，则存在更低维实现 $(\tilde{\mathbf{A}},\tilde{\mathbf{b}},\tilde{\mathbf{c}})$ 给出同一 $W(s)$；因 $\det(s\mathbf{I}-\tilde{\mathbf{A}})$ 阶次更低，$W(s)$ 必有分子分母公因子相消（即零极点对消）。\emph{充分性}：若无对消，则 $W(s)$ 已是“最简分式”，对应实现不可再降维，即最小实现，由 \cref{thm:co-minreal} 既可控又可观\cite{liubao2006modern}。
\end{proof}

\begin{pitfall}[对消不告诉你“坏在哪”——同一对消的三副面孔]
零极点对消\emph{表明}系统不可控\emph{或}不可观（或两者），但\textbf{不能区分到底坏在哪一侧}\cite{liubao2006modern}。以 $W(s)=\dfrac{s+2.5}{(s+2.5)(s-1)}$ 为例，分子分母公因子 $(s+2.5)$ 对消，可由\emph{三种}本质不同的实现产生：
\[
  \underbrace{\dot{\mathbf{x}}=\begin{pmatrix}1&0\\0&-2.5\end{pmatrix}\mathbf{x}+\begin{pmatrix}1\\1\end{pmatrix}u,\ \mathbf{y}=(1,0)\mathbf{x}}_{\text{(a) 可控但不可观}}
\]
\[
  \underbrace{\cdots+\begin{pmatrix}1\\0\end{pmatrix}u,\ \mathbf{y}=(1,1)\mathbf{x}}_{\text{(b) 不可控但可观}}\qquad
  \underbrace{\cdots+\begin{pmatrix}1\\0\end{pmatrix}u,\ \mathbf{y}=(1,0)\mathbf{x}}_{\text{(c) 既不可控又不可观}}
\]
三者传递函数相同（都被对消成 $\frac{1}{s-1}$），内部结构却迥异。\cref{fig:co-cancel} 把三种结构并排画出。\textbf{致命之处}：被对消掉的极点 $s=-2.5$（此例恰稳定）若换成\emph{不稳定}极点，则那个“消失”的模态会在系统内部悄悄发散、最终拖垮整个系统——这正是经典控制中“用零点抵消不稳定极点”这一危险做法的本质：它破坏了可控性 / 可观性，把不稳定性藏进了传递函数看不见的角落。\textbf{记住：极点不会因为对消而消失，它只是变得不可控或不可观了。}
\end{pitfall}

\begin{figure}[ht]
\centering
\begin{tikzpicture}[>=Stealth, font=\scriptsize,
  blk/.style={draw, rounded corners, minimum width=11mm, minimum height=7mm, align=center, fill=main!6, draw=main!60!black},
  iso/.style={draw, rounded corners, minimum width=11mm, minimum height=7mm, align=center, fill=third!10, draw=third!60!black, dashed},
  sum/.style={draw, circle, inner sep=0.5pt, minimum size=3.5mm},
  ln/.style={->, semithick, gray!55!black}]
\def\dx{4.9}
% (a)
\node at (0,1.9) {\textbf{(a) 可控·不可观}};
\node[blk] (a1) at (-0.7,0.9) {$\frac{1}{s-1}$};
\node[iso] (a2) at (-0.7,-0.5) {$\frac{1}{s+2.5}$};
\draw[ln] (-2.0,0.2) node[left]{$u$} |- (a1.west);
\draw[ln] (-2.0,0.2) |- (a2.west);
\node[sum] (ay) at (0.6,0.9) {$+$};
\draw[ln] (a1.east) -- (ay);
\draw[ln] (ay) -- ++(0.7,0) node[right]{$y$};
\node[third!60!black] at (0.55,-0.5) {不接 $y$};
% (b)
\node at (\dx,1.9) {\textbf{(b) 不可控·可观}};
\node[blk] (b1) at (\dx-0.7,0.9) {$\frac{1}{s-1}$};
\node[iso] (b2) at (\dx-0.7,-0.5) {$\frac{1}{s+2.5}$};
\draw[ln] (\dx-2.0,0.9) node[left]{$u$} -- (b1.west);
\node[third!60!black] at (\dx-1.55,-0.5) {不接 $u$};
\node[sum] (by) at (\dx+0.6,0.2) {$+$};
\draw[ln] (b1.east) -| (by);
\draw[ln] (b2.east) -| (by);
\draw[ln] (by) -- ++(0.7,0) node[right]{$y$};
% (c)
\node at (2*\dx,1.9) {\textbf{(c) 不可控·不可观}};
\node[blk] (c1) at (2*\dx-0.7,0.9) {$\frac{1}{s-1}$};
\node[iso] (c2) at (2*\dx-0.7,-0.5) {$\frac{1}{s+2.5}$};
\draw[ln] (2*\dx-2.0,0.9) node[left]{$u$} -- (c1.west);
\node[sum] (cy) at (2*\dx+0.6,0.9) {$+$};
\draw[ln] (c1.east) -- (cy);
\draw[ln] (cy) -- ++(0.7,0) node[right]{$y$};
\node[third!60!black] at (2*\dx-0.7,-1.15) {既不接 $u$ 也不接 $y$};
\end{tikzpicture}
\caption{同一传递函数 $\frac{s+2.5}{(s+2.5)(s-1)}$ 的三种实现：被对消的模态 $\frac{1}{s+2.5}$（虚框）分别因不接输出 (a)、不接输入 (b)、两端皆不接 (c) 而成为不可观 / 不可控 / 两者皆非。三者外部传递函数全等。（仿刘豹\cite{liubao2006modern} 图 3.15 重绘）}
\label{fig:co-cancel}
\end{figure}

\begin{note}
\textbf{多变量例外.}\;
\cref{thm:co-cancel} 仅对 SISO 严格成立。对 MIMO 系统，“传递函数阵无零极点对消”只是最小实现的\emph{充分}条件——即便出现某种对消，多变量系统仍\emph{可能}既可控又可观\cite{liubao2006modern}。MIMO 的可控可观与传递函数阵的关系要用 Smith--McMillan 形 / 不变因子来精确刻画，超出本基础章范围（用到时见 \cref{ch:robust-hinf} 的频域 MIMO 一节）。
\end{note}

% ════════════════════════════════════════════════════════════════
\section{离散与时变情形的补遗}
\label{sec:co-discrete}

\paragraph{离散系统：递推即判据。}
离散定常系统 $\mathbf{x}(k+1)=\mathbf{G}\mathbf{x}(k)+\mathbf{H}\mathbf{u}(k)$ 的可控性可由递推法\emph{直接看出}\cite{liubao2006modern}：把解 $\mathbf{x}(n)=\mathbf{G}^n\mathbf{x}(0)+\sum_{j=0}^{n-1}\mathbf{G}^{n-1-j}\mathbf{H}\mathbf{u}(j)$ 令 $\mathbf{x}(n)=\mathbf{0}$，得关于 $\{\mathbf{u}(j)\}$ 的线性方程组，其系数矩阵正是\textbf{离散可控性矩阵}
\begin{equation}
  \mathcal{C}=\begin{bmatrix}\mathbf{H}&\mathbf{G}\mathbf{H}&\cdots&\mathbf{G}^{n-1}\mathbf{H}\end{bmatrix},\qquad \text{可控}\iff\operatorname{rank}\mathcal{C}=n.
  \label{eq:co-disc-cmat}
\end{equation}
可观性矩阵 $\mathcal{O}=[\mathbf{C};\mathbf{C}\mathbf{G};\cdots;\mathbf{C}\mathbf{G}^{n-1}]$ 同理。形式与连续情形完全一致（$\mathbf{A}\to\mathbf{G},\mathbf{B}\to\mathbf{H}$），但\emph{推导更直白}——无须 Cayley--Hamilton 的级数收缩，递推 $n$ 步天然只出现 $\mathbf{G}^0,\dots,\mathbf{G}^{n-1}$。

\begin{note}
\textbf{OCR 勘误.}\;
教材原 OCR 稿中 \cref{eq:co-disc-cmat} 末列误作 $\dot{\mathbf{G}}^{n-1}\mathbf{H}$（多了导数点），正确应为 $\mathbf{G}^{n-1}\mathbf{H}$（离散系统无导数）。本书据递推结构取正确版。
\end{note}

\begin{insight}
\textbf{离散的两点差异。}\;
(1) \textbf{可控 $\neq$ 可达}：连续系统 $\mathrm{e}^{\mathbf{A}t}$ 恒可逆，故可控与可达等价（\cref{prop:co-reach}）；离散系统 $\mathbf{G}$ 可\emph{奇异}（如 $\mathbf{G}$ 有零特征值），此时“能驱到零态”（可控）与“能从零态驱到任意态”（可达）可能不一致——可控是较弱要求。(2) \textbf{多输入可少于 $n$ 步}：$n$ 阶系统若输入数 $r=n$ 且 $\mathbf{H}$ 非奇异，\emph{一步}即可把任意 $\mathbf{x}(0)$ 驱到原点（$\mathbf{H}\mathbf{u}(0)=-\mathbf{G}\mathbf{x}(0)$ 直接解出）\cite{liubao2006modern}。这对应导论 \cref{eq:ctrl-reach} 的离散视角。
\end{insight}

\paragraph{输出可控性（互校自现控02）。}
许多场合被控量是\emph{输出}而非状态，于是有\textbf{输出可控性}\cite{zhang2017modern}：系统 $\dot{\mathbf{x}}=\mathbf{A}\mathbf{x}+\mathbf{B}\mathbf{u},\ \mathbf{y}=\mathbf{C}\mathbf{x}$ 输出完全可控 $\iff$
\begin{equation}
  \operatorname{rank}\begin{bmatrix}\mathbf{C}\mathbf{B}&\mathbf{C}\mathbf{A}\mathbf{B}&\mathbf{C}\mathbf{A}^2\mathbf{B}&\cdots&\mathbf{C}\mathbf{A}^{n-1}\mathbf{B}\end{bmatrix}=q\quad(q=\dim\mathbf{y}).
  \label{eq:co-output-ctrl}
\end{equation}

\begin{pitfall}[状态可控性与输出可控性没有必然联系]
切勿想当然地以为“状态可控 $\Rightarrow$ 输出可控”或反之\cite{zhang2017modern}。反例：$\dot{\mathbf{x}}=\big(\begin{smallmatrix}-4&1\\2&-3\end{smallmatrix}\big)\mathbf{x}+\big(\begin{smallmatrix}1\\2\end{smallmatrix}\big)u,\ \mathbf{y}=(1,0)\mathbf{x}$。状态可控性矩阵 $[\mathbf{b}\ \mathbf{A}\mathbf{b}]=\big(\begin{smallmatrix}1&-2\\2&-4\end{smallmatrix}\big)$ 秩为 $1<2$——\emph{状态不可控}；而输出可控性矩阵 $[\mathbf{C}\mathbf{b}\ \mathbf{C}\mathbf{A}\mathbf{b}]=(1,-2)$ 秩为 $1=q$——\emph{输出可控}。二者分道扬镳。直觉：输出可控只要求“能操纵那 $q$ 个输出组合”，比“能操纵全部 $n$ 个状态”弱得多。
\end{pitfall}

% ════════════════════════════════════════════════════════════════
\section{常见误解汇总}
\label{sec:co-misconceptions}

\begin{enumerate}
  \item \textbf{“可控就是稳定 / 好控制”}——错。可控只说“能用输入操纵状态”，与稳定正交（\cref{sec:co-def} insight）。不稳定系统可完全可控（且正因可控才可被镇定）。
  \item \textbf{“对角型 $\mathbf{b}$ 无零行就一定可控”}——仅当特征值\emph{互异}或每个重特征值只一个约旦块时成立。两个约旦块同特征值时反例 $\mathbf{A}=\mathrm{diag}(1,1),\mathbf{b}=(1,1)^\top$ 不可控（\cref{sec:co-jordan} 陷阱）。
  \item \textbf{“多算几个 $\mathbf{A}^k\mathbf{b}$ 可能增秩”}——错。Cayley--Hamilton 保证算到 $\mathbf{A}^{n-1}$ 封顶，$\mathbf{A}^n$ 起不再贡献新方向（\cref{sec:co-kalman-cont} insight）。
  \item \textbf{“传递函数完整描述了系统”}——错。传递函数只反映可控且可观子块（\cref{thm:co-tf-co}）；不可控 / 不可观模态对它隐形，却可能（若不稳定）在内部发散。
  \item \textbf{“零极点对消把极点消掉了”}——错。极点不会消失，只是变得不可控或不可观（\cref{sec:co-cancel} 陷阱）；对消不可控 / 不可观侧无法仅凭传递函数区分。
  \item \textbf{“可控标准型转置就是可观标准型”}——仅 SISO 成立；MIMO 二者维数都可能不等（$nr$ vs $nm$，\cref{sec:co-canonical} 陷阱）。
  \item \textbf{“状态可控 $\Rightarrow$ 输出可控”}——错，二者无必然联系（\cref{sec:co-discrete} 陷阱）。
\end{enumerate}

% ════════════════════════════════════════════════════════════════
\section{通向高级：本章为谁打地基}
\label{sec:co-outlook}

本章建立的可控 / 可观、标准型、结构分解、最小实现，是 P4 控制部一整条上层主线的\emph{前提与 ground-truth}。沿单向交叉引用依赖链向上：

\begin{itemize}
  \item \textbf{综合（状态反馈 / 极点配置 / 观测器 / 分离原理）}——\cref{ch:linear-synthesis}：可控性是极点配置（Ackermann）的\emph{充要前提}，可观性是观测器存在的前提，\cref{eq:co-cstd} 的可控标准型 / \cref{eq:co-ostd} 的可观标准型是其天然设计坐标；可镇定 / 可检测（\cref{sec:co-stabilizable}）是输出反馈稳定的最弱要求。
  \item \textbf{LQR / LQG / Riccati}——\cref{ch:lqr-lqg-riccati}：CARE 存在唯一正定解的前提正是 $(\mathbf{A},\mathbf{B})$ \emph{可镇定} + $(\mathbf{A},\mathbf{Q}^{1/2})$ 可检测——本章把这两个条件讲透，C 章直接交叉引用、不再重述其内涵。
  \item \textbf{鲁棒 / 辨识 / 频域 $H_\infty$}——\cref{ch:robust-hinf}：子空间辨识恢复的就是某个\emph{最小实现}（\cref{thm:co-minreal}），其“维数唯一、坐标自由”（\cref{ex:co-minreal} insight）正是\emph{可辨识性}的根；MIMO 零极点对消（\cref{thm:co-cancel} note）的精确刻画也在该章。
  \item \textbf{安全控制 CLF/CBF}——\cref{ch:clf-cbf}：高相对度 CBF（ECBF）的反馈线性化用到 Brunovsky 标准型（可控标准型的多输入推广）与极点配置增益 $\mathbf{k}_b$——本章把标准型讲透，F 章直接交叉引用。
  \item \textbf{估计--控制对偶}——\cref{ch:eskf}：卡尔曼滤波的可观性正是本章 $(\mathbf{A},\mathbf{C})$ 的可观性，滤波收敛要求\emph{可检测}；LQR 增益与卡尔曼增益由\emph{对偶}（\cref{thm:co-duality}）的两个 Riccati 给出，这是“最优控制 $\leftrightarrow$ 最优估计”的桥。
\end{itemize}

\noindent 至此，可控可观这层“结构地基”已立。\textbf{下一章 \cref{ch:lyapunov-basics}}转向另一根支柱——稳定性与李雅普诺夫方法：本章问“能否操纵 / 反推”，下章问“放手后是否收敛”；二者合起来，才够撑起综合与最优控制的全部上层建筑。

% ════════════════════════════════════════════════════════════════
\section*{本章小结}

\begin{table}[H]\centering\small
\caption{符号表}
\begin{tabular}{lll}
\toprule
符号 & 含义 & 首次出现 \\
\midrule
$\mathbf{x},\mathbf{u},\mathbf{y}$ & 状态 / 输入 / 输出（维 $n,r,m$） & \cref{def:co-controllable} \\
$(\mathbf{A},\mathbf{B},\mathbf{C})$；$(\mathbf{G},\mathbf{H},\mathbf{C})$ & 连续 / 离散系统三元组 & \cref{def:co-controllable,eq:co-disc-cmat} \\
$\mathcal{C}=[\mathbf{B}\ \mathbf{A}\mathbf{B}\ \cdots\ \mathbf{A}^{n-1}\mathbf{B}]$ & 可控性矩阵（教材记 $\mathbf{M}$） & \cref{eq:co-cmat} \\
$\mathcal{O}=[\mathbf{C};\mathbf{C}\mathbf{A};\cdots;\mathbf{C}\mathbf{A}^{n-1}]$ & 可观性矩阵（教材记 $\mathbf{N}$） & \cref{eq:co-omat} \\
$\mathbf{W}_c,\mathbf{W}_o$ & 可控 / 可观 Gram 矩阵 & \cref{eq:co-tv-gramc} \\
$\boldsymbol{\Phi}(t,t_0)$ & 状态转移矩阵 & \cref{eq:co-control-relation} \\
$\mathbf{A}_c,\mathbf{b}_c$；$\mathbf{A}_o,\mathbf{c}_o$ & 可控 / 可观标准型矩阵 & \cref{eq:co-cstd,eq:co-ostd} \\
$\mathbf{R}_c,\mathbf{R}_o,\mathbf{R}$ & 结构分解变换阵 & \cref{eq:co-decomp-c,eq:co-kalman-4block} \\
$\bar{\mathbf{x}}_{co},\bar{\mathbf{x}}_{c\bar o},\bar{\mathbf{x}}_{\bar c o},\bar{\mathbf{x}}_{\bar c\bar o}$ & 四块状态（可控 / 可观各两态） & \cref{thm:co-kalman-decomp} \\
$W(s),\mathbf{W}(s)$ & 标量 / 矩阵传递函数 & \cref{eq:co-tf-co} \\
$\mathbf{D}=\lim_{s\to\infty}\mathbf{W}(s)$ & 直接传输项 & \cref{eq:co-D-limit} \\
\bottomrule
\end{tabular}
\end{table}

\begin{table}[H]\centering\small
\caption{判据与定理速查表}
\begin{tabular}{lll}
\toprule
判据 / 定理 & 一句话说明 & 对应 \\
\midrule
Kalman 秩判据 & 可控 $\iff\operatorname{rank}\mathcal{C}=n$（CH 压到 $n-1$ 项） & \cref{thm:co-rank} \\
约旦“一行 / 一列”法则 & 末行 / 首列非零（每特征值单块） & \cref{sec:co-jordan} \\
PBH 判据 & 逐模态：$\operatorname{rank}[\lambda\mathbf{I}-\mathbf{A}\ \mathbf{B}]=n$ & \cref{sec:co-stabilizable} \\
可镇定 / 可检测 & 只管不稳定模态；LQR/KF 的真前提 & \cref{sec:co-stabilizable} \\
时变 Gram 判据 & $\mathbf{W}_c$ 非奇异；秩判据的母判据 & \cref{thm:co-tv-gram} \\
对偶原理 & $(\mathbf{A},\mathbf{B})$ 可控 $\iff(\mathbf{A}^\top,\mathbf{B}^\top)$ 可观 & \cref{thm:co-duality} \\
可控 / 可观标准型 & 友矩阵；设计反馈 / 观测器的好坐标 & \cref{eq:co-cstd,eq:co-ostd} \\
Kalman 四块分解 & 可控可观 / 三块对传函隐形 & \cref{thm:co-kalman-decomp} \\
传函只见可控可观块 & $\mathbf{W}=\mathbf{C}_1(s\mathbf{I}-\mathbf{A}_{11})^{-1}\mathbf{B}_1$ & \cref{thm:co-tf-co} \\
最小实现 $\iff$ 可控且可观 & 维数唯一（McMillan 度），坐标自由 & \cref{thm:co-minreal} \\
SISO 无对消 $\iff$ 可控可观 & 极点不消失，只变不可控 / 不可观 & \cref{thm:co-cancel} \\
输出可控性 & $\operatorname{rank}[\mathbf{C}\mathbf{B}\ \cdots\ \mathbf{C}\mathbf{A}^{n-1}\mathbf{B}]=q$ & \cref{eq:co-output-ctrl} \\
\bottomrule
\end{tabular}
\end{table}

\begin{table}[H]\centering\small
\caption{知识点总表}
\begin{tabular}{clll}
\toprule
\# & 知识点 & 核心要点 & 难度 \\
\midrule
1 & 可控 / 可观定义 & 驱动能力 / 反映能力；可控$\neq$稳定 & \(\bigstar\bigstar\) \\
2 & 孤立方块直觉 & $b_i=0\Rightarrow$ 与输入脱钩的孤立模态 & \(\bigstar\bigstar\) \\
3 & Kalman 秩判据 & CH 把无穷步压到 $n$ 步 & \(\bigstar\bigstar\bigstar\) \\
4 & PBH / 可镇定可检测 & 逐模态；只管坏模态 & \(\bigstar\bigstar\bigstar\) \\
5 & 时变 Gram 判据 & 秩判据的母判据；LTV 刚需 & \(\bigstar\bigstar\bigstar\bigstar\) \\
6 & 对偶原理 & 一套结论用两次；物理意义相反 & \(\bigstar\bigstar\bigstar\) \\
7 & 可控 / 可观标准型 & 友矩阵；设计的天然坐标 & \(\bigstar\bigstar\bigstar\) \\
8 & Kalman 结构分解 & 四块；传函只见可控可观 & \(\bigstar\bigstar\bigstar\bigstar\) \\
9 & 最小实现 & 可控且可观；维数唯一坐标自由 & \(\bigstar\bigstar\bigstar\bigstar\) \\
10 & 零极点对消 (SISO) & 极点不消失只变不可控 / 不可观 & \(\bigstar\bigstar\bigstar\) \\
11 & 离散 / 输出可控性 & 可控$\neq$可达；状态$\neq$输出可控 & \(\bigstar\bigstar\bigstar\) \\
\bottomrule
\end{tabular}
\end{table}

\section*{延伸阅读}
\begin{itemize}
  \item \textbf{结构直觉与例题主线}：刘豹《现代控制理论》（第 3 版）第 3 章\cite{liubao2006modern}——本章孤立方块图像、Kalman 四块分解、最小实现、§3.10 零极点对消的主要出处。
  \item \textbf{标准形与最小实现互校}：张嗣瀛 / 高立群《现代控制理论》（第 2 版）第 4 章\cite{zhang2017modern}——最干净的友矩阵标准形、输出可控性、时变判据，作本章独立第二信源。
  \item \textbf{奠基原典}：Kalman《On the general theory of control systems》\cite{kalman1960general} 与《Mathematical description of linear dynamical systems》\cite{kalman1963mathematical}——可控可观与对偶原理的最初提出；卡尔曼滤波\cite{kalman1960new} 与之同源。
  \item \textbf{现代教材视角}：Chen《Linear System Theory and Design》\cite{chen1999linear}、Ogata《Modern Control Engineering》\cite{ogata2010modern}——PBH 判据、Smith--McMillan 形、MIMO 实现理论的标准参考。
\end{itemize}

\section*{本章与后续章节的关系}
\begin{table}[H]\centering\small
\begin{tabular}{lll}
\toprule
后续章节 & 关系 & 本章铺垫 \\
\midrule
\cref{ch:linear-synthesis} 综合 & 极点配置 / 观测器设计 & 可控可观前提、可控 / 可观标准型 \\
\cref{ch:lqr-lqg-riccati} LQR/LQG & CARE 唯一解前提 & 可镇定 + 可检测（\cref{sec:co-stabilizable}）\\
\cref{ch:robust-hinf} 辨识 / 鲁棒 & 子空间辨识恢复最小实现 & 最小实现、可辨识性（\cref{thm:co-minreal}）\\
\cref{ch:clf-cbf} CLF/CBF & ECBF 反馈线性化 & Brunovsky / 可控标准型、极点配置 \\
\cref{ch:eskf} 误差状态滤波 & 滤波可观性 = $(\mathbf{A},\mathbf{C})$ 可观 & 对偶原理、可观性矩阵 \\
\bottomrule
\end{tabular}
\end{table}

\section*{故障排查手册}
\begin{enumerate}
  \item \textbf{症状}：状态反馈极点\emph{配不动}某些极点 / Ackermann 求逆失败。\textbf{排查}：先验 $\operatorname{rank}\mathcal{C}=n$（\cref{thm:co-rank}）；不满秩则系统不可控、只能配可控子空间的极点（\cref{thm:co-decomp-c}）；若不可控的恰是不稳定模态，则系统\emph{不可镇定}（\cref{sec:co-stabilizable}），无解。
  \item \textbf{症状}：观测器误差\emph{不收敛}。\textbf{排查}：验 $\operatorname{rank}\mathcal{O}=n$（\cref{eq:co-omat}）；不满秩则不可观，至少要\emph{可检测}（不稳定模态可观）才有收敛观测器；用对偶（\cref{thm:co-duality}）把问题转成 $(\mathbf{A}^\top,\mathbf{C}^\top)$ 的可控 / 可镇定来查。
  \item \textbf{症状}：闭环传递函数“看起来稳”、实际系统\emph{内部发散}。\textbf{原因}：存在不可控 / 不可观的\emph{不稳定}模态，对传函隐形（\cref{thm:co-tf-co}）；常因设计时用零点对消了不稳定极点（\cref{sec:co-cancel} 陷阱）。\textbf{对策}：永远在状态空间核对全部 $n$ 个极点，勿信传函约分。
  \item \textbf{症状}：辨识 / 实现得到的模型\emph{维数偏高}、含冗余状态。\textbf{原因}：实现非最小（含不可控 / 不可观分量）。\textbf{对策}：做结构分解抽可控可观子块（\cref{thm:co-minreal}）；SISO 可先查传递函数有无零极点对消（\cref{thm:co-cancel}）。
  \item \textbf{症状}：沿轨迹线性化的 LTV 系统用常数秩判据\emph{判错}可控性。\textbf{原因}：LTV 不能用 $(\mathbf{A},\mathbf{B})$ 常数秩判据。\textbf{对策}：改用时变 Gram 判据（\cref{thm:co-tv-gram}）或其 $\mathbf{Q}_c(t)$ 充分判据（\cref{sec:co-timevarying} note），注意后者\emph{只充分}。
\end{enumerate}
